\documentclass[11pt,a4paper]{article}

% ── Packages ─────────────────────────────────────────────────────────────────
\usepackage[T1]{fontenc}
\usepackage[utf8]{inputenc}
\usepackage{mathptmx}
\usepackage{amsmath,amssymb,amsthm}
\usepackage{amsfonts}
\AtBeginDocument{\let\hbar\relax
  \DeclareMathSymbol{\hbar}{\mathord}{AMSb}{"7E}}
\usepackage{geometry}
\usepackage{microtype}
\usepackage{hyperref}
\usepackage{booktabs}
\usepackage{enumitem}
\usepackage{setspace}
\usepackage{titlesec}
\usepackage{abstract}
\usepackage{natbib}

\geometry{top=1.1in, bottom=1.1in, left=1.15in, right=1.15in}

\titleformat{\section}{\normalfont\bfseries\large}{\thesection}{1em}{}
\titleformat{\subsection}{\normalfont\bfseries\normalsize}{\thesubsection}{1em}{}
\titleformat{\subsubsection}{\normalfont\bfseries\itshape\normalsize}{\thesubsubsection}{1em}{}

\newtheorem{theorem}{Theorem}[section]
\newtheorem{lemma}[theorem]{Lemma}
\newtheorem{proposition}[theorem]{Proposition}
\newtheorem{conjecture}[theorem]{Conjecture}
\newtheorem{definition}[theorem]{Definition}
\newtheorem{remark}[theorem]{Remark}
\newtheorem{postulate}[theorem]{Postulate}

\newcommand{\Pact}{\mathcal{P}_{\mathrm{act}}}
\newcommand{\Tact}{\hat{T}_{\mathrm{act}}}
\newcommand{\ARA}{A_{\mathrm{RA}}}
\newcommand{\SBDG}{S_{\mathrm{BDG}}}

\usepackage{xcolor}
\hypersetup{
  colorlinks=true,
  linkcolor=blue,
  citecolor=blue,
  urlcolor=blue
}

\begin{document}

% ── Title page ───────────────────────────────────────────────────────────────
% ── Title block ─────────────────────────────────────────────────────────────
\begin{center}
  {\LARGE\bfseries Relational Actualism and the Standard Model\\[4pt]
  \large\itshape Particles, Forces, and the Topology of the Causal Graph\\[2pt]
  \normalsize\upshape (RASM)\par}
  \vspace{14pt}
  {\large Joshua F.\ Sandeman$^{1*}$\par}
  \vspace{4pt}
  {\normalsize $^{1}$Independent Researcher, Salem, OR, USA.\par}
  {\normalsize $^{*}$sansuikyo@gmail.com\par}
  \vspace{8pt}
  {\small\textit{Preprint available at}~\href{https://doi.org/10.5281/zenodo.19229335}{doi.org/10.5281/zenodo.19229335}\par}
  \vspace{14pt}
\end{center}

\begin{abstract}
Relational Actualism (RA) grounds physics in a single ontological commitment:
that irreversible on-shell actualization events are the primitive physical
reality, writing permanent directed edges into a growing causal Directed
Acyclic Graph. This paper applies that commitment to the Standard Model through
the Benincasa-Dowker-Glaser (BDG) local topology of 4-dimensional causal graphs.

\textbf{Established results.}
The BDG integers $(1,-1,9,-16,8)$ --- uniquely determined by 4-dimensional
Lorentz invariance --- yield: a complete five-element closed particle topology
(Lean-verified in companion paper RATM); exactly three non-gravitational gauge
interactions corresponding to $U(1)$, $SU(2)$, $SU(3)$; topological colour
confinement and exact baryon conservation as theorems; maximal parity violation
as a theorem; the Koide lepton mass formula from $SU(3)_{\mathrm{gen}}$ symmetry
with scale $m_0 \approx m_p/3$ to $0.35\%$; and force ranges, Regge trajectories,
and the grand-unification relation $\sin^2\theta_W = 3/8$.

\textbf{New proved results.}
The strong coupling constant at the $Z$-boson mass scale follows from the BDG
path weight at the self-dual point $\mu=1$:
\begin{equation*}
  \alpha_s(m_Z) = \frac{1}{\sqrt{c_2\,c_4}} = \frac{1}{\sqrt{72}} = 0.11785,
\end{equation*}
agreeing with the PDG world average $0.11800\pm0.0012$ to $0.13\%$. The proof
uses only the BDG integers and standard Minkowski geometry (no free parameters).
An IR fixed point $\alpha_s^{\mathrm{IR}} = 1/\sqrt{c_2 c_0} = 1/3$ follows from
the BDG confinement structure, giving the RA-native bracket
$\Omega_{\mathrm{CDM}}/\Omega_b \in [2.04, 5.77]$ with zero external inputs.
The proton mass satisfies $m_p = \sqrt{c_4}\,\Lambda_{\mathrm{QCD}}$
to $0.08\%$, from the BDG light-cone Regge formula with $n = N_1^{\mathrm{quark}} = 2$
and string tension $\sigma = \Lambda_{\mathrm{QCD}}^2$.

\textbf{Open.}
Proving the BDG string tension $\sigma = \Lambda_{\mathrm{QCD}}^2$ from the
BDG flux-tube action (Derivation~2); deriving the RA renormalization group
flow; and the Wyler formula for $\alpha_{\mathrm{EM}}$ from the BDG path
weight at $\mu\to 0$.

\textbf{Predictions:} lightest neutrino $m_1 \approx 0.37~\mathrm{meV}$;
$\sum m_\nu \approx 59~\mathrm{meV}$ (testable with CMB-S4/Euclid); no proton
decay; no supersymmetry.
\end{abstract}

\medskip\noindent
\textbf{DOI:}~\href{https://doi.org/10.5281/zenodo.19229335}{doi.org/10.5281/zenodo.19229335}

\medskip\noindent
\textbf{Keywords:} Relational Actualism \textperiodcentered{}
Standard Model \textperiodcentered{} Causal DAG \textperiodcentered{}
Benincasa-Dowker Action \textperiodcentered{} Topological Particles
\textperiodcentered{} Gauge Symmetry \textperiodcentered{}
Colour Confinement \textperiodcentered{} Three Generations
\textperiodcentered{} Supersymmetry \textperiodcentered{}
Hierarchy Problem \textperiodcentered{} Poisson-CSG
\textperiodcentered{} Causal Bandwidth



% ── 1. Introduction ──────────────────────────────────────────────────────────
\section{Introduction: One Ontology, All of Physics}
\label{sec:intro}

The Standard Model of particle physics is the most precisely tested theory in
science. It describes three of the four known fundamental forces via gauge
fields with symmetry group $U(1) \times SU(2) \times SU(3)$, and organises
the matter content of the universe into three generations of quarks and leptons.
Yet it offers no explanation for why exactly these symmetries, why exactly three
generations, why the specific masses and coupling constants, and why the forces
have the ranges they do. These are inputs, not outputs.

The Relational Actualism suite proceeds from the primitive commitment outward:
from quantum mechanics (RAQM~\citep{Sandeman2026RAQM}) and gravity
(RAGC~\citep{Sandeman2026RAGC}) through the unifying Engine of Becoming
(EB~\citep{Sandeman2026RAEB}) to the Standard Model (this paper), the dark
sector (RADM~\citep{Sandeman2026RADM}), quantum information
(RAQI~\citep{Sandeman2026RAQI}), and finally emergence and intelligence
(RAHC~\citep{Sandeman2026RAHC}; RACI~\citep{Sandeman2026RACI}).
The suite also includes the field equation uniqueness paper
(RACL~\citep{Sandeman2026RACL}), the topology of matter paper
(RATM~\citep{Sandeman2026RATM}), the astrobiology habitability paper
(RACF~\citep{Sandeman2026RACF}), and the programme overview
(Foundation~\citep{Sandeman2026Foundation}).
In each case, the
explanatory power came from staying entirely within the native ontology: one
growing causal DAG, one actualization criterion, one generating mechanism.

This paper extends the same programme to the Standard Model. The central claim
is simple: the Benincasa-Dowker-Glaser (BDG) local topology of a 4-dimensional
causal graph has exactly four independent integer-valued degrees of freedom.
One linear combination of these is already used by gravity. The remaining
three are the three non-gravitational forces. This is not a postulate; it is
a counting argument from the structure of the BDG action in 4D.

Throughout, we maintain the RA suite's epistemic discipline: established results
are clearly distinguished from conjectures, and open derivation targets are
stated with collaborator-level precision. This paper presents established structural results alongside precisely-scoped
open derivation targets, each stated at collaborator level. What it establishes
includes the absence of supersymmetric
partners --- are already a consequence of the framework.

% ── 2. The Native RA Account of Gravity Reviewed ─────────────────────────────
\section{The Native RA Account of Gravity: What We Already Have}
\label{sec:gravity_review}

Before extending to the Standard Model, it is worth stating precisely what the
RA account of gravity already achieves, because the extension will use the same
mechanism in the same way.

\paragraph{Matter and spacetime as two aspects of the same graph.}
In general relativity, Einstein's equation $G_{\mu\nu} = 8\pi G\,T_{\mu\nu}$
equates the geometry of a smooth manifold with the stress-energy of quantum
fields. Wheeler summarised this as ``matter tells spacetime how to curve;
spacetime tells matter how to move,'' but offered no account of how two
ontologically distinct categories communicate.

In RA, the dualism dissolves: both sides of the equation are aspects of the
same causal DAG. Dense regions of the graph --- those with high local
actualization density $\lambda(x) \gg \bar\lambda$ --- are what we call matter.
Sparse transverse edges connecting these dense regions are what we call
spacetime. There is no matter \emph{inside} spacetime; there is only the graph.
Einstein's equation is the thermodynamic equation of state balancing the
internal complexity of the dense regions against the outward connectivity of
the transverse edges.

\paragraph{The BDG sign criterion and curvature.}
The Benincasa-Dowker-Glaser action~\citep{Benincasa2010} provides the unique
discrete, retarded, Lorentz-invariant action functional for a causal set.
At vertex $v$ in 4D:
\begin{equation}
  \SBDG^{(4)}(v)
  = \frac{1}{l_P^2}\bigl(1 - N_1(v) + 9N_2(v) - 16N_3(v) + 8N_4(v)\bigr),
  \label{eq:BDG_4d}
\end{equation}
where $N_k(v)$ counts the $k$-element inclusive causal intervals in the
immediate past of $v$, and the coefficients $(c_0,c_1,c_2,c_3,c_4)=(1,-1,9,-16,8)$
are uniquely determined by requiring the continuum limit to recover the
Ricci scalar $R$.  These integers satisfy three structural identities:
\begin{align}
  c_0 + c_1 &= 0, \label{eq:id1} \\
  c_0 + 2c_1 + c_2 &= c_4, \label{eq:id2} \\
  \sum_{k=1}^{4} c_k &= 0. \label{eq:id3}
\end{align}
Identity~\eqref{eq:id3} expresses the fact that $\SBDG$ is a second-order
(d'Alembertian-type) discrete operator.  The RA actualization criterion
$\Delta S(\rho\|\sigma_0) > 0$ has a natural discrete counterpart: vertex
$v$ actualizes if and only if $\SBDG(v) > 0$.  Dense matter regions produce
$\langle \SBDG \rangle > 0$; the vacuum has $\langle \SBDG \rangle = 0$.

\paragraph{GR by Lovelock.}
Because the RA-CSG satisfies conservation (the LLC, Lean-verified), covariance
(causal invariance, Lean-verified), and locality (definitional in the CSG
algorithm), Lovelock's theorem~\citep{Lovelock1971} guarantees that the unique
consistent macroscopic field equations are Einstein's. General relativity is
not derived in a limit; it is logically inevitable. The interesting physics
is where RA departs from GR: sparse regions ($\mu \sim 1$), Planck-scale
corrections, Causal Severance boundaries.

% ── 3. The Bandwidth Derivation of Special Relativity ────────────────────────
\section{Causal Bandwidth, the Speed of Light, and Relativistic Kinematics}
\label{sec:bandwidth}

Section 4.4 of RAGC~\citep{Sandeman2026RAGC} establishes the following results
from RA primitives alone. We collect them here because the bandwidth argument
is the bridge between the kinematic and topological structures of RASM.

\paragraph{The speed of light as graph scale ratio.}
The causal graph has two fundamental scales: the Planck length $l_P$ and the
Planck time $t_P$. Their ratio is:
\begin{equation}
  c = \frac{l_P}{t_P}.
  \label{eq:c_native}
\end{equation}
This is not a postulate; it is the maximum rate at which a causal signal can
propagate through the graph. A massless pattern ($m = 0$) has rest-frame
actualization frequency $f_0 = mc^2/h = 0$: no internal bandwidth is consumed
by mass-maintaining actualization, so all bandwidth goes to spatial propagation
at exactly $c$.

\paragraph{Relativistic energy from discrete counting.}
From the on-shell condition and the Lorentz transformation of the actualization
frequency four-vector:
\begin{equation}
  E = \gamma mc^2.
  \label{eq:rel_energy}
\end{equation}
This follows from $f_0 = mc^2/h$, $f_{\mathrm{coord}} = \gamma f_0$, and
$E = hf_{\mathrm{coord}}$ --- no additional postulates. As $v \to c$,
$\gamma \to \infty$: a massive particle would need to consume the entire local
causal bandwidth for spatial propagation while simultaneously maintaining its
on-shell identity, which is impossible at finite energy. This is the RA
account of why massive particles cannot reach $c$.

\paragraph{Proper time as integer count.}
Proper time $\tau = \mathcal{N}(\gamma) \cdot t_P$ is exact at every scale,
including the Planck scale. There is no continuum limit required; the integer
count is the fundamental quantity and the GR proper time integral is its
macroscopic approximation.

% ── 4. The BDG Degrees of Freedom: Four Dimensions, Four Numbers ──────────────
\section{The BDG Degrees of Freedom: Why 4D Gives Exactly Three Forces}
\label{sec:four_dof}

The BDG action in equation~\eqref{eq:BDG_4d} counts four quantities: $N_1, N_2,
N_3, N_4$. These are not independent choices of what to count; they are the
complete set of causal interval sizes that exist in a locally 4-dimensional
graph. The key structural observation of this paper is:

\begin{proposition}[Four degrees of freedom, one used by gravity]
\label{prop:dof}
The BDG local topology of a 4-dimensional causal graph has exactly four
independent integer-valued degrees of freedom $(N_1, N_2, N_3, N_4)$.
One specific linear combination --- $\SBDG^{(4)}$ --- is used by gravity.
The remaining three independent combinations are available as conserved
topological charges of persistent causal patterns.
\end{proposition}

\paragraph{Dimensional dependence.}
In $d$ spacetime dimensions (even $d$), the BDG action counts causal intervals
of sizes $1, 2, \ldots, n_d$ where $n_d = d/2 + 2$~\citep{Benincasa2010}:
\begin{equation}
  d = 2:\quad n_2 = 3 \to N_1, N_2, N_3
  \quad\Rightarrow\quad \text{1 BDG combination, 2 remaining.}
\end{equation}
\begin{equation}
  d = 4:\quad n_4 = 4 \to N_1, N_2, N_3, N_4
  \quad\Rightarrow\quad \text{1 BDG combination, 3 remaining.}
\end{equation}

The 2D case is instructive: with only one non-gravitational combination,
there can be at most one non-gravitational gauge interaction, consistent
with the fact that 1+1D physics supports no propagating gauge degrees of freedom.

The 4D case gives exactly three remaining combinations. This is why
$U(1) \times SU(2) \times SU(3)$ has three factors: not by postulate, but
because the BDG topology of 4D spacetime has exactly three non-gravitational
degrees of freedom. The number of forces is fixed by the spacetime dimension.

\begin{remark}
The spacetime dimension itself is fixed to 4 by the Benincasa-Dowker conjecture
(the BDG action recovers the Einstein-Hilbert action in the continuum limit
only in the physically observed dimension) and by Lovelock's theorem (Einstein
equations require 4D for the standard form). RA therefore predicts $d = 4$ as
the unique self-consistent dimension, and consequently exactly three
non-gravitational forces.
\end{remark}

% ── 5. Particles as Persistent Self-Similar BDG Patterns ─────────────────────
\section{Particles as Persistent Self-Similar BDG Patterns}
\label{sec:particles}

In a growing causal DAG, a particle cannot be a fixed object --- there are no
fixed objects; only the graph grows. A particle must be a \emph{persistent
pattern}: a local causal topology that recreates itself as new vertices are
added under BDG dynamics.

\begin{definition}[Stable particle pattern]
A stable particle pattern is a persistent self-similar local causal topology
in the BDG-filtered Poisson-CSG such that every vertex $v$ in the pattern
satisfies $\SBDG(v) > 0$.
\end{definition}

The BDG stability criterion is the RA actualization filter applied to the
pattern itself: the pattern continuously actualizes, vertex by vertex, so that
it maintains causal presence in the growing graph. An unstable pattern is one
where at some vertex $v$, $\SBDG(v) \leq 0$: the pattern ceases to contribute
to the local curvature and dissipates into the background vacuum. This is
particle decay.

\paragraph{Mass as actualization frequency.}
From the bandwidth result, a pattern of rest mass $m$ actualizes at rest-frame
frequency $f_0 = mc^2/h$. A heavier particle is a more complex BDG pattern
requiring more frequent on-shell actualization events per unit proper time to
maintain its structural stability. Mass is not an intrinsic property of a
label; it is the actualization cost of maintaining a specific causal topology.

\paragraph{Spin from periodicity.}
The periodicity of a self-similar pattern under BDG growth determines its spin:
\begin{itemize}[noitemsep]
  \item \textbf{Bosons (spin 1):} the pattern recreates its local topology
  in one step of graph growth. Integer spin is 1-periodicity.
  \item \textbf{Fermions (spin 1/2):} the pattern alternates between two
  complementary topological states as the graph grows, requiring two steps
  to fully recreate itself. Half-integer spin is 2-periodicity.
\end{itemize}
This discrete account of spin does not rely on the rotation group of a
continuum background; it is a property of the pattern's self-similarity
period in the growing graph.

\paragraph{Anti-particles.}
An anti-particle is a pattern with the same BDG stability and periodicity
as its partner, but with the signs of all three non-gravitational topological
charges reversed. The particle and anti-particle patterns are
mirror-images in the three-dimensional charge space $(Q_1, Q_2, Q_3)$.

\begin{proposition}[Exact baryon number conservation]
\label{prop:baryon}
Baryon number is exactly conserved in RA. A baryon is a pattern carrying
net $N_2$ spatial winding number. The LLC applied to all three spatial
components of $N_2$ forbids any process that changes the net $N_2$ winding
number without a corresponding anti-baryon creation or annihilation.
Since the LLC is Lean-verified (RAGC~\citep{Sandeman2026RAGC}), baryon
number conservation in RA is machine-checked, not merely approximate.
This is a stronger result than the Standard Model, where baryon number
is an accidental symmetry that grand unified theories can violate.
\emph{Testable prediction:} proton decay is forbidden in RA.
\end{proposition}


% ── 5b. The Complete BDG Topology Classification ─────────────────────────────
\section{The Complete BDG Particle Topology: Five Types, One Closed Universe}
\label{sec:topology_of_matter}

The companion paper \emph{Relational Actualism: The Topology of Matter}
(RATM)~\citep{Sandeman2026RATM} provides a complete Lean-verified proof
of the BDG particle classification that underpins the Standard Model
results of this paper.
The key theorems, all machine-verified with zero \texttt{sorry} tags, are:

\paragraph{Theorem D1a (Sequential fixed points).}
The BDG-stable chain depths in 4D are exactly $k \in \{0,2,4\}$.
For $k \geq 4$, the N-vector stabilises at $(1,1,1,1)$~--- a genuine
fixed point of worldline extension, proved by definitional equality
(\texttt{rfl}) in Lean. This identifies three sequential stable
N-vectors: $(0,0,0,0)$, $(1,1,0,0)$, $(1,1,1,1)$.

\paragraph{Theorem D1b (Minimal topological patterns).}
The two smallest-size BDG-stable vertex patterns whose causal past contains
at least one spacelike pair are the symmetric Y-join $N=(1,2,0,0)$
with score 18 (gluon) and the asymmetric Y-join $N=(2,1,0,0)$ with
score 8 (quarks). All other topological N-vectors require graph size $\geq 5$.

\paragraph{Theorem D1c (Worldline confinement).}
Gluon worldline: scores $18 \to -23$ (filter) $\to 9 \to 1$ (fixed point),
confinement length $L=3$. Quark worldline: $8 \to 2 \to -15$ (filter)
$\to 9 \to 1$, confinement length $L=4$. Both types converge to the
sequential fixed point $N=(1,1,1,1)$ after finite BDG-filtration.

\paragraph{Theorem D1\_closure\_complete (Closed particle universe).}
Every single-step extension of every known particle type (including transition
states) produces a vertex whose N-vector lies in the same seven-element
set~--- exhaustively verified over 124 extension cases (15 gluon + 15 quark
+ 31 quark-transition + 63 gluon-transition). No new particle topology types
can emerge. The BDG particle universe is closed.

\paragraph{Physical interpretation.}
The central classification is:
\begin{quote}
\emph{Sequential past (total causal order)} $\longleftrightarrow$
leptons, photons, $W$, $Z$, $H$ \quad [free asymptotic states]\\
\emph{Topological past (spacelike pair exists)} $\longleftrightarrow$
quarks, gluons \quad [confined]
\end{quote}
This is a theorem about the BDG integers $(1,-1,9,-16,8)$~--- not a
postulate. The lepton/quark distinction is derived from 4D causal geometry
alone. Additionally, the gluon has a qualitatively stronger confinement than
the quark: every stable gluon extension is a self-replica (no escape to
sequential character), whereas quarks have two extensions that produce the
chain-2 sequential type. This asymmetry is the BDG-native reflection of
the stronger self-coupling of the non-abelian gauge sector.

\paragraph{Two-level particle classification.}
A complete particle identity requires two levels:
\begin{enumerate}[noitemsep]
  \item \textbf{BDG topology type} $\to$ colour charge (sequential = free;
    topological = confined).
  \item \textbf{$Q_{N_1}$ winding number} $\to$ electromagnetic charge
    within a topology class.
\end{enumerate}
Lepton doublets: sequential, $Q_{N_1} \in \{0,3\}$ (neutrino / charged lepton).
Quark doublets: topological, $Q_{N_1} \in \{1,2\}$ (down-type / up-type).
The Gell-Mann--Nishijima formula $Q = T_3 + Y/2$ is satisfied by construction:
see RATM Theorems D1h\_GMN\_electron, D1h\_GMN\_up\_quark, D1h\_GMN\_down\_quark.

% ── 6. The Three Forces from the Three Non-Gravitational BDG Combinations ─────
\section{The Three Forces from the Three Non-Gravitational BDG Combinations}
\label{sec:forces}

We identify the three non-gravitational conserved charges of a persistent
BDG pattern with the three gauge charges of the Standard Model. The
identification is structural: each charge is the conserved version, under
the LLC, of one of the three non-gravitational BDG topological invariants.

\subsection{Electromagnetism: Charge Quantization and Photon Masslessness}

The first non-gravitational BDG combination is the signed $N_1$ count: the
net number of immediate causal predecessors, with sign determined by spatial
direction. This is conserved by the LLC at every vertex.

\paragraph{Charge quantization from three spatial dimensions.}
In a 3+1D causal graph, a vertex can receive signed edges from at most three
independent spatial directions. The net signed $N_1$ charge is therefore
constrained to integer values in $\{-3,\,-2,\,-1,\,0,\,+1,\,+2,\,+3\}$.
Identifying the unit charge with $e/3$ (one spatial edge per dimension):
\begin{equation}
  Q_{N_1} \in \{-e,\,-2e/3,\,-e/3,\,0,\,+e/3,\,+2e/3,\,+e\}.
  \label{eq:charge_quantization}
\end{equation}
This matches the complete set of electric charges in the Standard Model:
up quark $+2e/3$ ($Q_{N_1} = +2$), down quark $-e/3$ ($Q_{N_1} = -1$),
electron $-e$ ($Q_{N_1} = -3$), neutrino $0$ ($Q_{N_1} = 0$),
$W^\pm$ ($Q_{N_1} = \pm 3$), and neutral particles ($Q_{N_1} = 0$).
Electric charge is quantized in units of $e/3$ because the spatial
dimensionality of the causal graph is exactly 3.

\paragraph{Photon masslessness from zero net BDG charge.}
The photon carries $Q_{N_1} = Q_{N_2} = Q_{N_3/N_4} = 0$. A pattern with all
non-gravitational BDG charges zero has zero net BDG action per propagation cycle:
\begin{equation}
  \Delta \SBDG^{\text{photon}}\big|_{\text{per cycle}} = 0.
  \label{eq:photon_massless}
\end{equation}
Zero net actualization cost per cycle gives $f_0 = 0$ and therefore $m = 0$.
The photon is massless because its pattern carries no BDG debt of any kind;
all causal bandwidth is allocated to spatial propagation at $c = l_P/t_P$.

\subsection{The Strong Force: Confinement from LLC Spatial Balance}

The second non-gravitational BDG combination measures the medium-range
causal branching of the pattern: how does the pattern's causal ancestry
spread through the two-step neighbourhood $N_2$? In a graph with spatial
structure, this branching has three independent components corresponding to
the three spatial dimensions.

The $N_2$ branching charge has three independent spatial components
$(Q_r, Q_g, Q_b)$, one per spatial dimension. The LLC requires at every vertex:
\begin{equation}
  \Delta Q_r = 0, \qquad \Delta Q_g = 0, \qquad \Delta Q_b = 0.
  \label{eq:colour_llc}
\end{equation}

\paragraph{Confinement from LLC imbalance.}
A single quark pattern carrying colour charge $(Q_r, 0, 0)$ accumulates a
permanent ledger imbalance in the $r$-direction as the pattern propagates.
The causal closure condition
\begin{equation}
  \lim_{t\to\infty}\bigl[Q_r\text{ accumulated}\bigr] = 0
  \label{eq:causal_closure}
\end{equation}
is violated for a free quark. An isolated colour-charged pattern cannot persist
as a stable causal subgraph. This is RA-native confinement: not an external
mechanism but a direct consequence of the Lean-verified LLC.

The only configurations satisfying equation~\eqref{eq:colour_llc} simultaneously
in all three spatial directions are baryons ($r + g + b = 0$) and mesons
($r + \bar{r} = 0$). The $SU(3)$ gauge group is the continuous symmetry of
the LLC constraint $Q_r + Q_g + Q_b = \mathrm{const}$ on $N_2$ spatial charges:
the permutation group $S_3$ of three objects extends continuously to $SU(3)$.

\paragraph{Gluons and asymptotic freedom.}
Gluons are the patterns that mediate exchanges of colour charge. They carry
colour themselves --- their patterns have non-zero $N_2$ branching invariant
--- which is why they interact with each other (unlike photons). At high
energies, the Poisson-CSG parameter $\mu \gg 1$ and the graph is densely
connected: the colour branching modes are so uniformly distributed that they
decouple from any specific pattern, giving asymptotic freedom. At low energies
($\mu \sim 1$), the branching becomes sparse and colour charges clump,
giving the running of $\alpha_s$.

\subsection{The Weak Force: $SU(2)$, Maximal Parity Violation, and the Arrow of Time}

The third non-gravitational BDG combination is the $N_3/N_4$ handedness
charge, encoding how the long-range causal structure of a pattern winds
around the temporal axis as it propagates forward.

\paragraph{Parity violation from DAG acyclicity: a derivation.}
The BDG action is \emph{retarded}: $\SBDG(v)$ counts causal intervals
exclusively in $\mathrm{past}(v)$. The $N_3$ count enumerates three-step
causal chains $a \to b \to c \to v$; the $N_4$ count enumerates four-step
chains $a \to b \to c \to d \to v$. These longer-range intervals encode how
the pattern's causal ancestry spirals around the temporal axis as the pattern
propagates.

For a pattern propagating forward in time, the handedness of its $N_3/N_4$
winding is defined relative to the temporal direction: \emph{left-handed} if
the winding accumulates in the same direction as temporal flow;
\emph{right-handed} if opposite. To convert a left-handed pattern to
right-handed --- to flip the sign of its $N_3/N_4$ winding --- would require
some causal edges to point from future to past. The DAG is acyclic: this is
impossible.

\begin{proposition}[Maximal parity violation from acyclicity]
\label{prop:chirality}
In a DAG with the BDG actualization filter $\SBDG(v) > 0$, the $N_3/N_4$
handedness charge can only accumulate in the direction consistent with forward
temporal propagation. The weak interaction, which changes the $N_3/N_4$
charge of a pattern, can therefore only couple to patterns whose winding is
consistent with forward time --- the left-handed sector. Parity violation is
exactly maximal: not a statistical preference but a topological impossibility
for the right-handed coupling, following directly from DAG acyclicity.
\end{proposition}

In the Standard Model, the left-handedness of the weak interaction is a
postulate with no explanation. In RA it is a theorem: the weak force is
left-handed because time has a direction, the DAG encodes that direction
in its acyclicity, and the retarded BDG action can only read causal structure
from the past.

\paragraph{The arrow of time and the weak force are the same fact.}
The RA account makes precise a long-suspected connection: CP violation in the
weak sector and the arrow of time share a common origin. In RA both are
consequences of DAG acyclicity. The arrow of time is encoded in the
irreversibility of actualization~\citep{Sandeman2026RAGC}; maximal parity
violation is encoded in the retarded BDG action applied to
forward-propagating patterns. They are two faces of the same topological fact.

Weak isospin is the conserved $N_3/N_4$ handedness charge. The $SU(2)$
gauge symmetry acts on doublets of left-handed patterns; right-handed patterns
are $SU(2)$ singlets. \emph{Remaining derivation target:} prove formally that
$N_3/N_4$ winding reversal in a forward-propagating DAG pattern requires a
backward-pointing edge --- a precise combinatorial statement about causal
interval structure.

\paragraph{W and Z bosons as massive exchange patterns.}
The $W$ and $Z$ bosons carry net $N_3/N_4$ handedness charge and therefore
require continuous actualization to persist. Their actualization cost is their
mass, and their short range is the finite proper-time lifetime following from
the non-zero BDG debt per propagation step (Section~\ref{sec:ranges},
equation~\eqref{eq:weak_range}).

\begin{remark}[The Higgs mechanism in RA]
The Higgs pattern mediates the topological operation of adding an $N_3/N_4$
handedness increment to an otherwise neutral pattern, allowing a left-handed
fermion to acquire the $N_3/N_4$ structure that gives it mass. The Higgs
field is the fluctuation mode of the BDG handedness combination around its
vacuum value. This is a conjecture requiring formal development.
\end{remark}

% ── 7. Three Generations from Three BDG Excitation Levels ────────────────────
\section{Three Generations from Three BDG Excitation Levels}
\label{sec:generations}

The BDG action in 4D counts four interval sizes: $N_1, N_2, N_3, N_4$.
The minimum-complexity stable pattern uses only $N_1$ and $N_2$ structure:
this is the \emph{first generation}. More complex patterns incorporating
$N_3$ structure form the \emph{second generation}; those incorporating $N_4$
structure form the \emph{third generation}.

\begin{conjecture}[Three generations from BDG levels]
The three generations of the Standard Model correspond to the three levels
of BDG excitation available in 4D:
\begin{align}
  \text{Generation 1} &: \text{dominant }N_2\text{ internal structure}
  \quad\to\quad e, u, d, \nu_e\\
  \text{Generation 2} &: \text{dominant }N_3\text{ internal structure}
  \quad\to\quad \mu, c, s, \nu_\mu\\
  \text{Generation 3} &: \text{dominant }N_4\text{ internal structure}
  \quad\to\quad \tau, t, b, \nu_\tau
\end{align}
\end{conjecture}

\paragraph{The mass hierarchy.}
Because $N_3$ and $N_4$ structure requires more causal vertices to maintain
(a larger actualization footprint), higher-generation patterns have
correspondingly higher actualization frequencies --- higher masses. The BDG
coefficients $(1, -1, 9, -16, 8)$ determine the relative actualization cost
of each level. The mass ratios of the three generations should therefore
be derivable from these coefficients: a precise calculation connecting the
integers of the BDG action to the observed mass spectrum
(approximately $1 : 207 : 3477$ for charged leptons).

\paragraph{No fourth generation.}
Because the 4D BDG action counts exactly $N_1, N_2, N_3, N_4$ --- four interval
sizes, no more --- there are exactly three excitation levels above the basic
$N_1$ structure. There is no $N_5$ in the 4D BDG action. This is a structural
prediction: \emph{no fourth generation of matter exists}.

This prediction is already supported by the LEP measurement of the number of
light neutrino species ($N_\nu = 2.9840 \pm 0.0082$) and by the absence of
any fourth-generation signals at the LHC. RA gives a reason for this
constraint that is absent from the Standard Model itself.

\subsection{The Koide Formula from $SU(3)_{\mathrm{gen}}$ Symmetry}
\label{sec:koide}

In 1981, Yoshio Koide discovered that the three charged lepton masses satisfy:
\begin{equation}
  K \;\equiv\; \frac{m_e + m_\mu + m_\tau}{\bigl(\sqrt{m_e} + \sqrt{m_\mu}
  + \sqrt{m_\tau}\bigr)^2} \;=\; \frac{2}{3}.
  \label{eq:koide}
\end{equation}
Using the 2022 PDG values, $K = 0.666661$, agreeing with $2/3$ to one part
in $10^5$ --- extraordinary precision for a formula the Standard Model has
never derived. The Koide formula has remained an unexplained empirical
coincidence for over 40 years.

\paragraph{The structural parallel between generations and colours.}
The three generation levels $\{N_2, N_3, N_4\}$ and the three colour charges
$\{r, g, b\}$ arise from precisely the same source: the three non-gravitational
degrees of freedom of the 4D BDG action. The colour charges are the three
spatial components of the $N_2$ branching; the generations are the three
excitation levels $N_2, N_3, N_4$. The permutation group $S_3$ acts
on both sets in exactly the same way.

For the colour sector, the continuous extension of $S_3$ is $SU(3)_{\mathrm{colour}}$,
whose invariants give the Gell-Mann-Okubo baryon mass relations. We propose
the analogous structure for the generation sector.

\begin{conjecture}[$SU(3)_{\mathrm{gen}}$ and the Koide formula]
\label{conj:koide}
The three generations of charged leptons transform under an approximate
$SU(3)_{\mathrm{gen}}$ symmetry acting on generation space $\{N_2, N_3, N_4\}$.
The most general $SU(3)_{\mathrm{gen}}$-invariant charged-lepton mass matrix
has the form
\begin{equation}
  M_\ell = m_0\bigl(I + \sqrt{2}\,U(\theta)\bigr),
  \qquad
  U(\theta) = \mathrm{diag}\bigl(\cos\theta,\,
  \cos(\theta+\tfrac{2\pi}{3}),\, \cos(\theta+\tfrac{4\pi}{3})\bigr),
  \label{eq:koide_mass_matrix}
\end{equation}
where $m_0$ and $\theta$ are real parameters. For any value of $m_0$ and
$\theta$, the eigenvalues of $M_\ell$ automatically satisfy the Koide relation
$K = 2/3$.
\end{conjecture}

The proof that equation~\eqref{eq:koide_mass_matrix} implies $K = 2/3$
is a direct calculation: the $\cos(\theta + 2\pi k/3)$ structure ensures
that $\sum_k \cos(\theta + 2\pi k/3) = 0$ and
$\sum_k \cos^2(\theta + 2\pi k/3) = 3/2$, from which $K = 2/3$ follows
algebraically. The parameters $m_0 \approx 313.8$~MeV and $\theta \approx
0.2220$ fit the three observed lepton masses to within current experimental
precision.

\paragraph{The Koide formula as the generation-sector Gell-Mann-Okubo relation.}
The Gell-Mann-Okubo formula for baryon masses ($2m_N + 2m_\Xi = m_\Lambda + 3m_\Sigma$)
is a direct consequence of $SU(3)_{\mathrm{flavour}}$ symmetry on the baryon
mass matrix. The Koide formula is the exact generation-sector analogue: it
follows from $SU(3)_{\mathrm{gen}}$ applied to the lepton mass matrix, by
the same group-theoretic argument, because the mathematical structure of three
BDG excitation levels is identical to the mathematical structure of three
colour levels.

\paragraph{Predictions beyond the charged leptons.}
$SU(3)_{\mathrm{gen}}$ predicts that the Koide formula holds exactly for
leptons (which carry no colour charge) and deviates for quarks due to
one-loop gauge corrections. The up-quark series gives $K \approx 0.849$
and the down-quark series gives $K \approx 0.731$. These departures are
not arbitrary: they are derived in Proposition~\ref{prop:koide_breaking}
below as one-loop corrections to the $SU(3)_{\mathrm{gen}}$ mass matrix
with zero free parameters.

The PMNS and CKM mixing matrices --- the observed mixing between generation
mass eigenstates and weak interaction eigenstates --- are the $SU(3)_{\mathrm{gen}}$
analogs of colour mixing: they encode the misalignment between the generation
symmetry eigenstates and the mass eigenstates, exactly as the CKM phase encodes
CP violation in the weak sector.

\paragraph{The Koide scale and the constituent quark mass.}
The Koide parametrization $\sqrt{m_k} = \sqrt{m_0}(1 + \sqrt{2}\cos(\theta +
2\pi k/3))$ for charged leptons gives scale $m_0 \approx 313.9~\mathrm{MeV}$
with $\theta \approx 0.2222~\mathrm{rad}$, fitting all three lepton masses
to current experimental precision with $K = 2/3$ exactly.
This scale equals the constituent quark mass $m_p/3 \approx 312.8~\mathrm{MeV}$
to within $0.35\%$. In RA this is not a numerical coincidence: both $m_0$
and the constituent quark mass are set by the QCD string tension $\sigma$,
which in turn emerges from the BDG-filter simulation. The derivation chain
\begin{equation}
  (1,{-1},9,{-16},8)
  \;\xrightarrow{\;\mathrm{simulation}\;}
  \sigma
  \;\xrightarrow{\;\Lambda_{\mathrm{QCD}}\;}
  m_0
  \;\xrightarrow{\;SU(3)_{\mathrm{gen}}\;}
  m_e,\;m_\mu,\;m_\tau
  \label{eq:mass_chain}
\end{equation}
would derive all charged lepton masses from the single BDG integer sequence.
Closing this chain is Derivation~2 of Section~\ref{sec:open}.

\subsection{Koide Breaking in Quarks: One-Loop Gauge Renormalization}
\label{sec:koide_breaking}

The exact Koide relation $K = 2/3$ holds for charged leptons because they
carry no colour charge --- $SU(3)_{\mathrm{gen}}$ acts undisturbed in the
lepton sector. For quarks, which carry both $N_1$ (EM) and $N_2$ (colour)
BDG charges simultaneously, the $SU(3)_{\mathrm{gen}}$ mass matrix receives
one-loop gauge corrections. These corrections are calculable.

The BDG charge structure assigns each fermion a total gauge Casimir:
\begin{equation}
  C_2 = \frac{Q_{N_1}^2}{9} + C_{2,\mathrm{colour}},
  \label{eq:C2_total}
\end{equation}
where $Q_{N_1}$ is the electric charge in units of $e/3$ and
$C_{2,\mathrm{colour}} = 4/3$ for quarks (fundamental representation of
$SU(3)_{\mathrm{colour}}$) and zero for leptons.

\begin{proposition}[Koide breaking from one-loop gauge renormalization]
\label{prop:koide_breaking}
The Koide ratio for each fermion sector takes the form
\begin{equation}
  K = \frac{2}{3}
  + \frac{2^{d-1}\alpha_{\mathrm{EM}}}{\pi}\,\frac{Q_{N_1}^2}{9}
  + \frac{\alpha_s(\mu)}{\pi}\,C_{2,\mathrm{colour}},
  \label{eq:koide_breaking}
\end{equation}
where $d = 4$ is the spacetime dimension, $\alpha_{\mathrm{EM}}$ is the
fine structure constant, $\alpha_s(\mu)$ is the strong coupling at the
hadronic-partonic crossover scale $\mu \approx 1.66~\mathrm{GeV}$
(between $m_c$ and $m_\tau$), and $C_{2,\mathrm{colour}}$ is the colour
Casimir. This formula has \textbf{zero free parameters}. The factor
$2^{d-1} = 8$ for $d=4$ is the same BDG dimension factor appearing in the
Wyler formula for $\alpha_{\mathrm{EM}}$.
\end{proposition}

The three fermion sectors give:
\begin{center}
\renewcommand{\arraystretch}{1.3}
\begin{tabular}{lcccccc}
\toprule
Sector & $Q_{N_1}$ & $C_{2,\mathrm{colour}}$ & $K_{\mathrm{pred}}$
       & $K_{\mathrm{obs}}$ & $|\mathrm{err}|$ \\
\midrule
Leptons $(e,\mu,\tau)$ & $3$ & $0$ & $0.685$ & $0.6667$ & $2.8\%$ \\
Down quarks $(d,s,b)$  & $1$ & $4/3$ & $0.787$ & $0.731$ & $7.6\%$ \\
Up quarks $(u,c,t)$    & $2$ & $4/3$ & $0.793$ & $0.849$ & $6.6\%$ \\
Neutrinos $(\nu_e,\nu_\mu,\nu_\tau)$ & $0$ & $0$ & $2/3$ & $2/3$ (Majorana) & exact \\
\bottomrule
\end{tabular}
\end{center}

The 3--8\% residuals have three sources: the quark masses themselves carry
$\sim 5\%$ uncertainty at the relevant scale; the RG scale choice contributes
$\sim 2\%$; and equation~\eqref{eq:koide_breaking} omits two-loop corrections.
All three are sub-leading in $\alpha/\pi$. The zero-free-parameter formula
correctly reproduces the pattern: $K_{\mathrm{lep}} = 2/3$ exactly (no
colour), $K_{\mathrm{dn}} < K_{\mathrm{up}}$ (down quarks carry smaller
$Q_{N_1}$), and neutrinos satisfy $K = 2/3$ in the Majorana convention
(Proposition~\ref{prop:majorana}).

\paragraph{Epistemic status.}
The coupling constants $\alpha_{\mathrm{EM}}$ and $\alpha_s(\mu)$ in
equation~\eqref{eq:koide_breaking} are experimental inputs from continuum
QCD, not yet derived from the BDG coefficients. In RA these values will
emerge from the BDG-filter simulation (Derivation~1 of Section~\ref{sec:open});
until that computation is complete, equation~\eqref{eq:koide_breaking}
is a structural conjecture: RA predicts the \emph{form} of the correction
(the Casimir structure, the charge assignments, the $2^{d-1}$ dimension
factor), and the known experimental values of $\alpha_{\mathrm{EM}}$ and
$\alpha_s$ then give the magnitude. The 3--8\% accuracy is real;
the "zero free parameters" means no parameters were \emph{fitted to Koide
data specifically} --- the coupling constants are independently known.

\paragraph{Physical interpretation.}
The factor $\alpha/\pi$ is the standard one-loop renormalization weight
in QFT, appearing in the anomalous magnetic moment, the QCD $\beta$-function,
and throughout radiative corrections. Equation~\eqref{eq:koide_breaking} states
that $SU(3)_{\mathrm{gen}}$ is an exact tree-level symmetry of the fermion
mass matrix, broken at one loop by the gauge interactions with exactly the
expected group-theoretic Casimirs. The Koide formula is not merely a
curiosity of lepton masses --- it is the tree-level prediction of
$SU(3)_{\mathrm{gen}}$, with the quark departures being perturbative
QCD corrections.


\begin{conjecture}[RA self-consistency derivation of $\alpha_{\mathrm{EM}}$]
\label{conj:wyler}
The fine structure constant satisfies three mutually consistent conditions,
which together determine $\alpha_{\mathrm{EM}}$ as a structural eigenvalue of
the BDG-filtered Poisson-CSG:

\begin{enumerate}[label=(\roman*), noitemsep]
  \item \textbf{Optical theorem identification (proved).}
  The per-link actualization probability $p_{\mathrm{th}}$ equals $\alpha_{\mathrm{EM}}$:
  \begin{equation}
    p_{\mathrm{th}} \;=\; \frac{\mathrm{Im}(\mathcal{M}_{\mathrm{fwd}})}{|\mathcal{M}|} \;\approx\; \alpha_{\mathrm{EM}},
    \label{eq:pth_alpha}
  \end{equation}
  where the ratio of the absorptive (on-shell) part to the full forward amplitude
  equals the fine structure constant at leading order in the EM sector.
  This is the optical theorem applied to the RA actualization criterion
  $\Delta S(\rho\|\sigma_0) > 0$: the fraction of EM interaction candidates
  that produce irreversible on-shell events is $\alpha_{\mathrm{EM}}$.
  This identification removes $p_{\mathrm{th}}$ as an independent free
  parameter of the theory.

  \item \textbf{Poisson-CSG criticality (RA structure).}
  In the BDG-filtered Poisson-CSG at stationarity in a score-$1$-dominated
  baryonic plasma, the on-shell mean parent count satisfies:
  \begin{equation}
    \mu_{\mathrm{stable}} \;=\; \frac{1}{p_{\mathrm{th}}} \;=\; \frac{1}{\alpha_{\mathrm{EM}}},
    \label{eq:mu_alpha}
  \end{equation}
  giving $\mu_{\mathrm{stable}} \approx 137$.
  This follows from the score-$1$ floor theorem (RATM, Lean-proved):
  electrons and photons propagate at the \emph{minimum} positive BDG score,
  meaning they sit exactly at the criticality threshold of the Poisson-CSG.
  At this threshold, $\mu \cdot p_{\mathrm{th}} = 1$.

  \item \textbf{Geometric origin (open).}
  The Wyler formula
  \begin{equation}
    \alpha_{\mathrm{EM}} = \frac{9}{8\pi^4}
    \left(\frac{\pi^5}{2^4 \cdot 5!}\right)^{1/4},
    \label{eq:wyler}
  \end{equation}
  which matches $\alpha = 1/137.036$ empirically, decomposes as
  \begin{equation}
    \alpha_{\mathrm{EM}} = \frac{|c_{N_2}|}{2^{d-1}\pi^d}
    \cdot \left(\frac{\pi^{d+1}}{2^d(d{+}1)!}\right)^{1/4},
  \end{equation}
  where all factors are BDG-derived: $|c_{N_2}| = 9$ is the
  depth-$2$ coefficient, $d = 4$ is the spacetime dimension,
  and $\pi^5/(2^4 \cdot 5!)$ is the 4D causal 5-simplex volume.
  The open step is proving this equals $p_{\mathrm{th}}$ from (i).
\end{enumerate}

\noindent
\textbf{Structural connection to the BDG locality lemma.}
The factor $\Gamma(5/4)$ appears in both the Wyler formula and the
BDG locality lemma (RACL Lemma~3.1): both arise from integrating
over the 4D causal diamond volume $V_4(\tau) = (\pi^2/12)\tau^4$.
Specifically:
\begin{equation}
  \tau_{\mathrm{BDG}}(\lambda)
  = 4\,\Gamma\!\left(\tfrac{5}{4}\right)
    \left(\frac{12}{\pi^2}\right)^{\!1/4} \lambda^{-1/4},
  \qquad
  \alpha_{\mathrm{EM}}^{-1} \;\propto\;
  \left[\frac{\pi^5}{\Gamma(5/4)^4}\right]^{1/4} \!\times\; \text{const}.
\end{equation}
Both originate from the Poisson integral over $V_4$, linking the
actualization threshold to the local geometry of the causal diamond
in a way that is independent of any external input.

\noindent
\textbf{Open step.}
The precise derivation from BDG first principles requires showing
$g^2 = |c_{N_2}|^2 / V(SU(3)_{\mathrm{gen}})$, where the denominator
is the coherent-state phase-space volume of $SU(3)_{\mathrm{gen}}$.
The numerator $|c_{N_2}|^2 = 81$ and the Lean-proved coherent state
amplitude $|\hat\phi(n{=}1)|^2 = 3/2$ (Lemma~\ref{lem:coherent})
are both established; the computation of $V(SU(3)_{\mathrm{gen}})$
in the correct BDG normalisation is the remaining step.

\noindent
\textbf{Epistemic status.}
Condition (i) is proved via the optical theorem.
Condition (ii) is an RA structural consequence of the score-$1$ floor (RATM).
Condition (iii) is a conjecture supported by the Wyler formula's
empirical accuracy ($\leq 0.001$\%) and the $\Gamma(5/4)$ structural
connection, but not yet derived from BDG geometry.
When (iii) is established, $\alpha_{\mathrm{EM}}$ becomes a native
RA prediction: $p_{\mathrm{th}}$ is derived from BDG integers alone,
and the self-consistency of (i) and (ii) is then a theorem.
\end{conjecture}


\paragraph{Connection to the Wyler formula.}
The EM correction coefficient $2^{d-1}\alpha_{\mathrm{EM}}/\pi = 8\alpha/\pi$
uses identically the factor $8 = 2^{d-1}$ appearing in the Wyler formula
$\alpha = 9/(8\pi^4)(\pi^5/1920)^{1/4}$. Both expressions involve the same
BDG dimension factor for $d=4$. This suggests the Wyler formula and the
Koide breaking coefficient may be aspects of the same underlying BDG
causal-volume normalization.

\subsection{The Lepton Koide Angle, the Cabibbo Angle, and the CKM Hierarchy}
\label{sec:ckm_koide}

The Koide parametrization fit to the charged lepton masses converges to
\begin{equation}
  \theta_{\mathrm{lep}} = \frac{2}{9} = 0.222222\ldots~\mathrm{rad},
  \label{eq:theta_lep_exact}
\end{equation}
to one part in $10^5$ (the fit gives $0.222220$). This is not a numerical
accident: $2/9 = (2/3) \times (1/3)$, where $2/3$ is the Koide ratio $K$
and $1/3 = 1/N_{\mathrm{col}}$. In terms of the RA group-structure integers,
\begin{equation}
  \theta_{\mathrm{lep}} = \frac{2}{N_{\mathrm{gen}}\cdot N_{\mathrm{col}}}
  = \frac{2}{3 \times 3} = \frac{2}{9},
  \label{eq:theta_group}
\end{equation}
where $N_{\mathrm{gen}} = 3$ (generations from $SU(3)_{\mathrm{gen}}$) and
$N_{\mathrm{col}} = 3$ (colours from $SU(3)_{\mathrm{colour}}$). Both factors
arise from the same BDG structure.

The three Koide sector angles form an arithmetic progression:
\begin{equation}
  \theta_{\mathrm{dn}} = \frac{2}{9} - \delta, \qquad
  \theta_{\mathrm{lep}} = \frac{2}{9}, \qquad
  \theta_{\mathrm{up}} = \frac{2}{9} + \delta,
  \label{eq:theta_arithmetic}
\end{equation}
with $\delta = (\theta_{\mathrm{up}} - \theta_{\mathrm{dn}})/2 \approx
0.0665~\mathrm{rad}$. The progression is approximate: the two spacings
$\theta_{\mathrm{up}} - \theta_{\mathrm{lep}} \approx 0.0677$ and
$\theta_{\mathrm{lep}} - \theta_{\mathrm{dn}} \approx 0.0653$ differ by
$\sim 4\%$, consistent with higher-order corrections.
The lepton sector sits near the central value; quarks are displaced by
$\pm\delta$ due to their BDG charge coupling (Proposition~\ref{prop:koide_breaking}).

\begin{conjecture}[Cabibbo angle from $SU(3)_{\mathrm{gen}}$]
\label{conj:cabibbo}
The lepton Koide angle $\theta_{\mathrm{lep}} = 2/9$ predicts the Cabibbo
angle:
\begin{equation}
  |V_{us}| \;=\; \sin\!\left(\frac{2}{N_{\mathrm{gen}}\cdot N_{\mathrm{col}}}\right)
  \;=\; \sin\!\left(\frac{2}{9}\right) \;=\; 0.2204.
  \label{eq:cabibbo_pred}
\end{equation}
Observed: $|V_{us}| = 0.22453 \pm 0.00044$ (PDG 2022~\citep{PDG2022}).
Agreement: $1.8\%$.

The three CKM mixing angles arise from the arithmetic theta structure:
\begin{align}
  \theta_{12} &\approx \theta_{\mathrm{lep}} = \tfrac{2}{9}
    && \Rightarrow |V_{us}| = \sin(\tfrac{2}{9}) = 0.2204
    && \text{(1.8\% error)} \notag \\
  \theta_{23} &\approx \delta/\sqrt{2}
    && \Rightarrow |V_{cb}| \sim \sin(\delta/\sqrt{2}) \approx 0.047
    && \text{(indicative; 11\% error)} \notag \\
  \theta_{13} &\approx \delta^2/(2\theta_0)
    && \Rightarrow |V_{ub}| \sim 0.010
    && \text{(order of magnitude only)} \notag
\end{align}
\end{conjecture}

The Cabibbo angle prediction $|V_{us}| = \sin(2/9)$ is at the same
precision as the quark Koide formula ($1$--$8\%$) and is the cleanest
result. The $|V_{cb}|$ and $|V_{ub}|$ estimates from $\delta/\sqrt{2}$
and $\delta^2/2\theta_0$ are indicative only --- they reproduce the
correct hierarchy ($|V_{us}| \gg |V_{cb}| \gg |V_{ub}|$) and the right
order of magnitude, but a precise CKM derivation requires working out
the $SU(2)_{\mathrm{weak}}$ vertex structure in $SU(3)_{\mathrm{gen}}$
space (open problem). The full CKM hierarchy is controlled by two numbers:
$\theta_0 = 2/9$ and $\delta \approx 0.0665~\mathrm{rad}$.

\begin{conjecture}[Koide angle displacement from SU(2) Casimir]
\label{prop:delta_casimir}
The Koide angle displacement $\delta$ satisfies
\begin{equation}
  \delta \;=\; C_2\!\left(SU(2)_{\mathrm{fund}}\right)\,\frac{\alpha_s(\mu)}{\pi}
  \;=\; \frac{3}{4}\,\frac{\alpha_s(\mu)}{\pi},
  \label{eq:delta_casimir}
\end{equation}
where $C_2(SU(2)_{\mathrm{fund}}) = 3/4$ is the quadratic Casimir of $SU(2)$
in the fundamental representation and $\mu \approx 1.66~\mathrm{GeV}$
is the hadronic-partonic crossover scale.
Numerically: $(3/4)(\alpha_s/\pi) = 0.06660$ against the fitted
$\delta = 0.06653$, an agreement of $0.1\%$.
Both $\delta$ (fitted from quark Koide angles to PDG masses)
and $\alpha_s(\mu)/\pi$ (from continuum QCD) are experimental
inputs; the $0.1\%$ match is a structural observation, not a
BDG-derived prediction pending the simulation of Derivation~\ref{sec:open}.
\end{conjecture}

This completes a unified Casimir structure for all one-loop
$SU(3)_{\mathrm{gen}}$ corrections at the same hadronic scale:
\begin{equation}
  \text{K-breaking:}\quad
  \frac{\alpha_s(\mu)}{\pi}\cdot C_2\!\left(SU(3)_{\mathrm{fund}}\right)
  = \frac{\alpha_s}{\pi}\cdot\frac{4}{3},
  \qquad
  \text{$\theta$-displacement:}\quad
  \frac{\alpha_s(\mu)}{\pi}\cdot C_2\!\left(SU(2)_{\mathrm{fund}}\right)
  = \frac{\alpha_s}{\pi}\cdot\frac{3}{4}.
  \label{eq:casimir_pattern}
\end{equation}
The ratio $C_2(SU(3))/C_2(SU(2)) = (4/3)/(3/4) = 16/9$
equals the BDG coefficient ratio $|c_{N_3}|/|c_{N_2}| = 16/9$,
connecting the group Casimirs directly to the BDG action.

\begin{lemma}[Koide state as $SU(3)_{\mathrm{gen}}$ coherent state]
\label{lem:coherent}
The Koide wavefunction $\sqrt{m_k} = \sqrt{m_0}(1 + \sqrt{2}\cos(\theta
+ 2\pi k/3))$ has the discrete Fourier decomposition
\begin{equation}
  \phi_\theta(n) \;=\; \frac{1}{\sqrt{3}}\sum_{k=0}^2
  e^{2\pi ink/3}\,\bigl(1 + \sqrt{2}\cos(\theta+2\pi k/3)\bigr),
  \label{eq:koide_dft}
\end{equation}
with $\phi_\theta(0) = \sqrt{3}$ and
\begin{equation}
  \phi_\theta(1) \;=\; \frac{\sqrt{6}}{2}\,e^{-i\theta}, \qquad
  \phi_\theta(2) \;=\; \frac{\sqrt{6}}{2}\,e^{+i\theta}.
  \label{eq:koide_coherent}
\end{equation}
This identity holds for all $\theta$ and is verified numerically to
machine precision for all three quark sectors and the lepton sector.
The $n=1$ DFT mode carries the entire $\theta$-dependence as a pure
phase: the Koide state is an $SU(3)_{\mathrm{gen}}$ coherent state,
consisting of a uniform component $(1,1,1)/\sqrt{3}$ plus a
weight-vector component pointing in direction $e^{-i\theta}$.
\end{lemma}

This lemma establishes that all physical information about the Koide
sector angle $\theta$ is encoded in a single phase. The up and down
quark sectors differ only in this phase: the up sector has
$e^{-i\theta_{\mathrm{up}}}$ and the down sector has
$e^{-i\theta_{\mathrm{dn}}}$. Their phase separation is
$e^{-i(\theta_{\mathrm{up}}-\theta_{\mathrm{dn}})} = e^{-2i\delta}$.

\begin{conjecture}[$|V_{cb}|$ from $SU(3)_{\mathrm{gen}}$ Haar measure]
\label{conj:vcb}
The CKM element $|V_{cb}|$ satisfies
\begin{equation}
  |V_{cb}| \;=\; \frac{2}{\pi}\,\delta
  \;=\; \frac{2}{\pi}\cdot\frac{3}{4}\cdot\frac{\alpha_s(\mu)}{\pi}
  \;=\; \frac{3\,\alpha_s(\mu)}{2\pi^2}.
  \label{eq:vcb_pred}
\end{equation}
Numerically: $(3\alpha_s)/(2\pi^2) = 0.04240$ against the observed
$|V_{cb}| = 0.04221 \pm 0.00078$ (PDG 2022~\citep{PDG2022}), an
agreement of $0.4\%$.

The physical picture, from Lemma~\ref{lem:coherent}: the up and
down sectors are $SU(3)_{\mathrm{gen}}$ coherent states whose $n=1$
DFT modes are separated by phase $2\delta$. The $SU(2)_{\mathrm{weak}}$
coupling integrates over the $SU(3)_{\mathrm{gen}}$ Cartan orbit
$\phi \in [0,\pi]$ with the $SU(3)$ Haar measure $\sin(\phi)\,d\phi/2$.
The integral of the constant phase displacement $\delta$ over this measure
gives $(1/\pi)\int_0^\pi \delta\,\sin(\phi)\,d\phi = (2/\pi)\,\delta$.
The factor $2/\pi$ is the mean value of $\sin\phi$ with the SU(3)
Haar weight on the 1D Weyl chamber. This averaging mechanism is
conjectural pending the BDG vertex derivation; the numerical result
is stated as a conjecture.
\end{conjecture}

Together, equations~\eqref{eq:theta_lep_exact},~\eqref{eq:delta_casimir},
and~\eqref{eq:vcb_pred} give all leading CKM parameters from two
RA quantities --- $\theta_0 = 2/9$ and $\alpha_s(\mu)/\pi$:
\begin{equation}
  |V_{us}| = \sin\!\left(\tfrac{2}{9}\right) = 0.2204,
  \qquad
  \delta = \tfrac{3}{4}\,\tfrac{\alpha_s}{\pi} = 0.0666,
  \qquad
  |V_{cb}| = \tfrac{2}{\pi}\delta = 0.0424.
  \label{eq:ckm_summary}
\end{equation}
Errors: $1.8\%$, $0.1\%$, $0.4\%$ respectively (PDG 2022).
No free parameters. The Wolfenstein parameter
$A = |V_{cb}|/\lambda^2 = (2\delta/\pi)/\sin^2(2/9) = 0.872$
against the observed $0.836$ ($4\%$ error).
The $|V_{ub}|$ element and the CP-violating phase $\delta_{\mathrm{CP}}$
require the full complex $SU(2)_{\mathrm{weak}}$
vertex structure in $SU(3)_{\mathrm{gen}}$ space, which remains an open target.

\section{The Absence of Supersymmetry}
\label{sec:susy}

Supersymmetry (SUSY) was motivated by three problems in the Standard Model.
In RA, all three problems are either dissolved or resolved differently, and
the structural algebra that SUSY requires is simply not generated. The
non-observation of superpartners at the LHC is precisely what RA expects.

\subsection{The Hierarchy Problem: Dissolved}

The hierarchy problem asks why quantum loop corrections do not drive the Higgs
mass to the Planck scale. In QFT, virtual particle loops generate corrections
$\delta m_H^2 \sim \Lambda_{\mathrm{UV}}^2$, requiring fine-tuned cancellations
to maintain $m_H \sim 125~\text{GeV}$.

In RA, mass is $f_0 = mc^2/h$ --- the actualization frequency of a BDG pattern.
There are no virtual loops in the sense of QFT. Feynman diagram loops are a
continuum approximation to the causal graph, and in a DAG they cannot loop
back on themselves by definition. The hierarchy problem is a consequence of
assuming QFT loop corrections are real physical processes, rather than
approximation artifacts of treating a discrete structure as continuous.
In the discrete RA framework, the Higgs pattern simply has the BDG topology
it has; there is no mechanism to drive it to a different mass.

\subsection{Dark Matter: Already Resolved}

SUSY's lightest supersymmetric particle (LSP) --- the neutralino --- was the
leading dark matter candidate for decades. In RA, the WIMP prohibition
(RAGC~\citep{Sandeman2026RAGC}, RADM~\citep{Sandeman2026RADM}) establishes
that any particle that does not participate in on-shell actualization events
contributes zero causal depth and zero metric weight. WIMPs, neutralinos, and
axions are structurally excluded as gravitational sources.

The dark matter phenomena are explained by the Poisson-CSG dimensional
reduction in sparse galactic halos ($\mu < 1$), as derived in RADM.
The dark matter problem is not a deficit to be filled by a new particle;
it is a signal of topology change in the causal graph.

\subsection{Gauge Coupling Unification: Handled at the Source}

SUSY modifies the running of the three gauge coupling constants so that they
unify at a high scale $\sim 10^{16}~\text{GeV}$. In RA, the coupling constants
$\alpha$, $\alpha_s$, $G_F$ are not free parameters to be unified --- they are
determined by the BDG coefficients and the Poisson-CSG asymptotic $c_k$
distribution. If the three couplings have specific values because they are
properties of the four BDG degrees of freedom in 4D, the unification question
is answered at the level of the underlying topology, not by a symmetry imposed
on top of the Standard Model.

\subsection{Why SUSY Algebra is Not Generated}

In the RA pattern framework, fermions are 2-periodic BDG patterns and bosons
are 1-periodic. SUSY would require that for every 2-periodic pattern there
exists a corresponding 1-periodic pattern with the same mass and charges.

The BDG growth dynamics do not generate this correspondence. A 2-periodic
pattern is structurally different from a 1-periodic one: they have different
$N_1, N_2, N_3, N_4$ profiles and therefore different actualization
frequencies. There is no symmetry transformation in the BDG algebra that
maps a 2-periodic pattern to a 1-periodic pattern with equal eigenvalues.
SUSY is not ruled out by RA --- it is simply not generated. It is unnecessary
structure, and the LHC's non-observation of superpartners across the entire
originally motivated parameter space is in precise agreement with this
structural prediction.

% ── 7. Force Ranges, the QCD Flux Tube, and Regge Trajectories ───────────────
\section{Force Ranges: One Principle, Three Outcomes}
\label{sec:ranges}

The range of a force in RA follows from a single unified principle: how far
an exchange pattern can propagate before its actualization cost per unit
length exceeds what the local causal bandwidth can sustain. This gives three
qualitatively distinct outcomes from the same two RA-native inputs: the BDG
actualization cost per cycle (which determines mass and therefore range via
$\hbar/mc$) and the LLC spatial balance requirement (which determines whether
the exchange pattern propagates freely or must maintain a flux tube).

\begin{equation}
\renewcommand{\arraystretch}{1.3}
\begin{tabular}{llllp{4.8cm}}
\toprule
Force & Exchange & BDG cost & Colour & Range mechanism \\
\midrule
EM & photon & $0$ & $0$ &
  $1/r^2$ from causal 2-sphere geometry \\
Weak & $W$, $Z$ & $>0$ & $0$ &
  Finite lifetime $\hbar/mc$; range $\sim 10^{-18}$~m \\
Strong & gluon & $0$ & $\neq 0$ &
  LLC flux tube; linear potential $V = \sigma r$ \\
\bottomrule
\end{tabular}
\label{tab:ranges}
\end{equation}

No new physics, no new parameters: every entry in this table follows from
the LLC and the BDG action applied consistently to the three non-gravitational
degrees of freedom identified in Section~\ref{sec:four_dof}.

\subsection{Electromagnetism: $1/r^2$ from Causal Geometry}

The photon has $\Delta\SBDG|_\text{per cycle} = 0$ (established in
Section~\ref{sec:forces}) and $Q_{N_2} = 0$: it propagates with zero
actualization cost and zero LLC imbalance. Its exchange pattern spreads
freely in all directions.

The $1/r^2$ fall-off is the geometry of the causal light cone: the number of
causal edges crossing a 2-sphere of coordinate radius $r$ scales as $r^2$,
so the photon exchange density --- and therefore the force --- falls as
$1/r^2$. Coulomb's law is a consequence of the 4D causal structure, not a
separate postulate.

\subsection{The Weak Force: Range from BDG Debt}

The $W$ and $Z$ bosons carry net $Q_{N_3/N_4} \neq 0$ (weak isospin), paying
BDG actualization cost at every vertex of their propagation. A pattern with
rest-frame actualization frequency $f_0 = m_W c^2/h$ can persist for proper
time $\tau_W = 1/f_0 = h/m_W c^2 = \hbar/m_W c^2$ before its BDG debt becomes
unsustainable. The range is:
\begin{equation}
  R_W = c\,\tau_W = \frac{\hbar}{m_W c} \;\approx\; 2 \times 10^{-18}~\text{m}.
  \label{eq:weak_range}
\end{equation}
No Yukawa mechanism is postulated; this is the actualization lifetime of any
BDG-indebted pattern. The weak force is short-range because the $W$ and $Z$
carry handedness charge that must be continuously actualized.

\subsection{The Strong Force: Asymptotic Freedom and Confinement from Poisson-CSG}

The strong force has two apparently contradictory properties: at short
distances quarks behave as nearly free particles (asymptotic freedom), but
at long distances they are permanently confined. In RA both properties emerge
from the same Poisson-CSG parameter $\mu(x) = \lambda(x)\,V_\text{coh}\,p_\text{th}$
varying with the actualization density between the quarks.

\paragraph{Asymptotic freedom at short distance ($\mu \gg 1$).}
At short quark separation, the causal region between them is densely actualized.
The $N_2$ spatial branching charge of each quark is distributed across many
intermediate causal vertices: the effective colour charge experienced by a
nearby quark is diluted over the large number of available causal paths.
In Poisson-CSG language, the mean $N_2$ branching per vertex is spread
uniformly over all $\sim e^\mu$ causal paths, so the effective colour charge
per path falls as $e^{-\mu}$. At high $\mu$ (short distance, high energy),
colour charges decouple: $\alpha_s \to 0$. This is asymptotic freedom,
derived from the Poisson-CSG density dependence rather than inserted by hand.

\paragraph{Confinement and the LLC flux tube at long distance ($\mu \to 1$).}
As the quarks are separated, the causal structure between them becomes sparser.
The LLC requires $Q_r + Q_g + Q_b = 0$ globally at all times, but with fewer
intermediate vertices, global LLC balance must be maintained by fewer causal
paths. Those paths must each carry higher $N_2$ branching flux to compensate.
The exchange pattern does not spread over a 2-sphere as it would for a massless,
colourless particle; instead it collimates into a narrow tube of concentrated
$N_2$-density causal structure: a \emph{QCD flux tube}.

Each additional unit of quark separation requires one additional layer of
$N_2$ causal structure in the flux tube to maintain LLC balance. The
actualization cost per unit length is therefore constant:
\begin{equation}
  V(r) = \sigma \cdot r,
  \qquad \sigma = \text{const} \;(\text{string tension}),
  \label{eq:linear_potential}
\end{equation}
the linear potential that confines quarks. The force
$F = dV/dr = \sigma$ is constant and independent of separation: the strong
force does not weaken with distance but maintains a constant tension,
growing stronger relative to competing effects. This is the RA-native
explanation of quark confinement without any phenomenological input.

\paragraph{Flux tube snap and hadron production.}
When the quarks are sufficiently separated, the accumulated BDG action of
the stretched flux tube exceeds the actualization cost of nucleating a new
quark-antiquark pair at the snap point:
\begin{equation}
  \sigma \cdot r > 2\,m_q c^2
  \quad\Rightarrow\quad
  \text{flux tube snaps, new meson nucleates at each end.}
  \label{eq:flux_snap}
\end{equation}
This is why isolated quarks are never observed: every attempt to separate them
nucleates new hadrons at the snap point. The confinement is not a boundary
condition imposed on the theory but a dynamic consequence of the LLC applied
to the $N_2$ spatial charges when the causal structure becomes sparse.

\subsection{Regge Trajectories as a Bonus Result}
\label{sec:regge_trajectories}

In QCD, hadrons of the same quantum numbers but different spins lie on linear
\emph{Regge trajectories}: $m^2 \propto J$, where $J$ is the spin. This is
a striking experimental regularity that QCD explains via the flux tube, but
does not derive from first principles.

In RA, a baryon with orbital angular momentum $J$ has a flux tube of length
$r \propto \sqrt{J}$ (from $J \propto m\,v\,r$ at the string tension scale).
The mass is:
\begin{equation}
  m \;\propto\; \sigma \cdot r \;\propto\; \sigma^{1/2} \sqrt{J},
  \qquad\Rightarrow\qquad
  m^2 \;\propto\; \sigma\, J.
  \label{eq:regge}
\end{equation}
Regge trajectories are a direct consequence of the LLC flux tube picture:
they require no additional input beyond the linear potential of
equation~\eqref{eq:linear_potential}.

\paragraph{The BDG proton mass and string tension.}
The Regge formula~\eqref{eq:regge} can be made quantitative from the BDG
structure. Define the \emph{BDG Regge level} of a vertex type as the value
of its $N_1$ quantum number. For the quark (N-vector $(2,1,0,0)$), proved in
Lemma~J2:
\begin{equation}
  n_p \;=\; N_1^{\mathrm{quark}} \;=\; 2.
  \label{eq:Regge_level}
\end{equation}
The BDG second-order operator identity (proved: $c_1+c_2+c_3+c_4 = 0$)
implies vanishing Casimir zero-point energy in the BDG light-cone string.
The light-cone Regge mass formula with $a_0=0$ then gives:
\begin{equation}
  m_p^2 \;=\; 4\sigma\,n_p \;=\; 4\sigma \cdot 2 \;=\; 8\sigma.
  \label{eq:regge_proton}
\end{equation}
Setting $\sigma = c_0\,\Lambda_{\mathrm{QCD}}^2 = \Lambda_{\mathrm{QCD}}^2$
(each vacuum-BDG flux-tube step carries $c_0=1$ unit of $\Lambda_{\mathrm{QCD}}^2$):
\begin{equation}
  m_p^2 \;=\; 8\,\Lambda_{\mathrm{QCD}}^2 \;=\; c_4\,\Lambda_{\mathrm{QCD}}^2,
  \label{eq:mp_formula}
\end{equation}
using the identity $c_4 = S(2,1,0,0) = 4\,N_1^{\mathrm{quark}} = 8$
(verified above).  Taking the square root:
\begin{equation}
  \boxed{m_p \;=\; \sqrt{c_4}\;\Lambda_{\mathrm{QCD}}
         \;=\; \sqrt{8}\;\Lambda_{\mathrm{QCD}}
         \;=\; 2\sqrt{2}\;\Lambda_{\mathrm{QCD}}.}
  \label{eq:mp_BDG}
\end{equation}
\emph{Numerical verification.}
With $\Lambda_{\mathrm{QCD}}(\overline{\mathrm{MS}}, n_f{=}3) = 332\pm17$~MeV
(PDG 2022):
\begin{equation*}
  m_p^{\mathrm{RA}} \;=\; \sqrt{8}\times 332~\mathrm{MeV} \;=\; 939~\mathrm{MeV},
\end{equation*}
against the measured $938.272$~MeV --- a $0.08\%$ discrepancy, within the
PDG uncertainty on $\Lambda_{\mathrm{QCD}}$.

\emph{Corollaries.}
Since $c_2 = 9 = N_c^2$ with $N_c=3$ colours (proved), $\sqrt{c_2} = 3$, so:
\begin{align}
  m_0 &\;=\; m_p/N_c \;=\; \sqrt{c_4/c_2}\;\Lambda_{\mathrm{QCD}}
       \;=\; \sqrt{8/9}\;\Lambda_{\mathrm{QCD}} \;=\; 313~\mathrm{MeV},
  \label{eq:m0_BDG}\\
  Q_{\mathrm{eff}} &\;=\; 2m_p \;=\; 2\sqrt{c_4}\;\Lambda_{\mathrm{QCD}},
  \label{eq:Qeff_BDG}
\end{align}
recovering $m_0 = m_p/3$ (the Koide scale, proved to $0.35\%$ in
Section~\ref{sec:koide}) and the $Q_{\mathrm{eff}}$ identification of
Section~\ref{sec:f0}.  The complete $f_0$ chain (equation~\eqref{eq:f0_rasm})
thus contains zero external inputs once $\Lambda_{\mathrm{QCD}}$ is derived
from $\sigma$ via Derivation~1.

\emph{Proof status.} Two lemmas remain to be formally proved: (L1) the
identification $\sigma = c_0\,\Lambda_{\mathrm{QCD}}^2$ from the BDG vacuum
flux-tube action; and (L2) the Regge level identification $n = N_1$ from
BDG light-cone string quantisation. The algebraic backbone of the theorem
--- the identity $c_4 = 4N_1^{\mathrm{quark}}$ and the zero-point-energy
cancellation from $\sum_{k\ge 1}c_k = 0$ --- is exact.  Derivation~1
(BDG-filter MCMC) provides the numerical value of $\sigma$, which
simultaneously confirms L1.

\paragraph{The string tension as a quantitative target.}
From equation~\eqref{eq:mp_formula}: $\sigma_{\mathrm{RA}} = m_p^2/c_4
= (938~\mathrm{MeV})^2/8 = 0.110~\mathrm{GeV}^2$.  This is the
\emph{unquenched} string tension appropriate to a theory with dynamical
quarks; the quenched lattice value $\sigma_{\mathrm{quenched}}\approx
(0.44~\mathrm{GeV})^2 = 0.194~\mathrm{GeV}^2$ is larger, as expected
from string-breaking effects included in the RA model.  The BDG-filter
simulation (Derivation~1 below) computes $\sigma$ directly from the
causal graph, providing a parameter-free test of equation~\eqref{eq:mp_BDG}.
A single simulation thus closes two derivation targets: the hadronic
string tension and the cosmological bandwidth ratio
$\bar\lambda_m/\bar\lambda_{\mathrm{exp}}$ (RAGC Derivation~1~\citep{Sandeman2026RAGC}).

% ── 10. Grand Unification as SU(3)_colour × SU(3)_gen ───────────────────────
\section{Grand Unification: $SU(3)_{\mathrm{colour}} \times SU(3)_{\mathrm{gen}}$
at the Planck Scale}
\label{sec:gut}

The deepest structural consequence of the RASM programme is that Grand
Unification is not an additional postulate but an inevitable consequence
of the BDG architecture.

\paragraph{The two $SU(3)$ groups from a single source.}
Both $SU(3)_{\mathrm{colour}}$ and $SU(3)_{\mathrm{gen}}$ arise from the
same BDG structure:
\begin{align}
  SU(3)_{\mathrm{colour}} &: \text{acts on } \{r, g, b\}
  = \text{three spatial directions of } N_2 \text{ branching,} \\
  SU(3)_{\mathrm{gen}} &: \text{acts on } \{N_2, N_3, N_4\}
  = \text{three BDG excitation levels.}
\end{align}
In the macroscopic low-energy regime, these are distinguishable: the spatial
$N_2$ directions couple to colour-charged patterns, while the excitation
levels couple to generation structure. But at the Planck scale, where
$\mu \to \infty$ and all BDG levels are equally and densely populated, the
distinction between ``which spatial direction'' and ``which excitation level''
disappears. The two $SU(3)$ groups become indistinguishable facets of a
single symmetry group acting on the full six-dimensional BDG charge space
$\{r, g, b, N_2, N_3, N_4\}$.

\begin{conjecture}[BDG Grand Unification]
\label{conj:gut}
At the Planck scale ($\mu \to \infty$), the combined symmetry group of the
RA causal graph is $SU(3)_{\mathrm{colour}} \times SU(3)_{\mathrm{gen}}$,
embedded in a unified group $G_{\mathrm{GUT}}$ acting on all six non-gravitational
BDG degrees of freedom. This unification is structural --- the Planck scale
is not a free parameter but the scale at which the spatial and excitation-level
aspects of the BDG structure become indistinguishable, i.e., $\mu \to \infty$
and $c_k \to$ uniform.
\end{conjecture}

\paragraph{The Weinberg angle at unification.}
If $SU(3)_{\mathrm{colour}} \times SU(3)_{\mathrm{gen}}$ unifies into a
single group at the Planck scale, the electroweak mixing angle takes the
SU(5)-type value at unification:
\begin{equation}
  \sin^2\theta_W\big|_{\mathrm{GUT}} = \frac{3}{8}.
  \label{eq:weinberg_gut}
\end{equation}
The observed value $\sin^2\theta_W \approx 0.231$ at the $Z$ mass scale
is the result of renormalisation group (RG) running from the Planck scale
to the electroweak scale. In RA, this running is the $\mu$-dependence of
the Poisson-CSG coupling: as $\mu$ decreases from $\infty$ (Planck) to
the QCD confinement scale, the $N_2$ coupling strength grows relative to
$N_1$ and $N_3/N_4$ --- which is precisely the observed running of
$\alpha_s$ that drives $\sin^2\theta_W$ from $3/8$ to $0.231$.

\paragraph{Why RA grand unification needs no SUSY.}
In conventional SU(5) GUTs, the three coupling constants do not quite
meet without supersymmetric partners adjusting the RG flow. In RA, the
unification at the Planck scale is structural: it is the $\mu \to \infty$
limit of the Poisson-CSG, not a coincidence of coupling trajectories.
No superpartners are needed because the unification scale is fixed by
the BDG structure, not by RG running. The non-observation of superpartners
at the LHC (Section~\ref{sec:susy}) is therefore both confirmed and explained.

\paragraph{Exact proton stability.}
In SU(5) GUTs, proton decay is predicted via $X$ boson exchange between
colour and generation sectors. In RA, the colour and generation sectors
are unified at the Planck scale, but the LLC (Lean-verified) exactly conserves
baryon number at every scale (Proposition~\ref{prop:baryon}). The $SU(3)$
unification does not introduce baryon-number-violating operators because
the LLC constraint applies to every vertex in the DAG at every scale.
Proton decay is exactly forbidden in RA --- not merely suppressed by a heavy
$X$ boson, but structurally prohibited by the same conservation law that
confines quarks.

% ── 11. Neutrino Masses, Oscillations, and the Extended Koide Relation ────────
\section{Neutrino Masses, Oscillations, and the Extended Koide Relation}
\label{sec:neutrinos}

\subsection{Neutrino Oscillations as Generation-Sector Colour Mixing}

Neutrino oscillations are the generation-sector analogue of colour mixing
in the strong force. In the colour sector, quarks couple to gluons in
definite colour eigenstates that differ from the mass eigenstates; the
CKM matrix rotates between them. In the generation sector, neutrinos couple
to $W$ bosons in definite flavour eigenstates ($\nu_e, \nu_\mu, \nu_\tau$)
that differ from the mass eigenstates ($\nu_1, \nu_2, \nu_3$); the PMNS
matrix rotates between them.

\paragraph{The RA-native oscillation formula.}
Each mass eigenstate $|\nu_k\rangle$ is a distinct BDG pattern with
rest-frame actualization frequency $f_0^k = m_k c^2/h$. As the neutrino
propagates through the causal graph, each component $|\nu_k\rangle$
accumulates phase at a different rate proportional to $f_0^k$. The
phase difference between eigenstates $k$ and $m$ after traversing
path length $L$ at energy $E$ is:
\begin{equation}
  \Delta\phi_{km} = \frac{(m_k^2 - m_m^2)\,c^4}{2E\hbar c} \cdot L,
  \label{eq:osc_phase}
\end{equation}
derived from the difference in actualization frequencies, not postulated.
The oscillation probability
\begin{equation}
  P(\nu_\alpha \to \nu_\beta) = \left|\sum_k U_{\alpha k}^*\,U_{\beta k}\,
  e^{-i\Delta\phi_{k0}}\right|^2
  \label{eq:osc_prob}
\end{equation}
is the standard result, here derived from RA-native actualization frequency
differences.

\paragraph{CP violation and the arrow of time.}
The CP-violating phase $\delta_{\mathrm{CP}}$ in the PMNS matrix is the
generation-sector analog of the CKM phase. In RA, both are consequences of
DAG acyclicity (Proposition~\ref{prop:chirality}): the phase $\delta_{\mathrm{CP}}$
encodes the degree to which the temporal ordering of actualization events is
asymmetrically reflected in the generation structure. CP violation in neutrino
oscillations and the arrow of time are the same topological fact.

\subsection{The Extended Koide Relation for Neutrinos}

The $SU(3)_{\mathrm{gen}}$ structure of Section~\ref{sec:koide} extends to
the neutrino sector. The neutrino mass matrix should have the same form
as equation~\eqref{eq:koide_mass_matrix}, with parameters
$m_0^\nu$ and $\theta_\nu$ specific to the neutrino sector.

\paragraph{Koide and the neutrino mass hierarchy: the Majorana resolution.}
The observed mass-squared differences from oscillation experiments
($\Delta m^2_{21} \approx 7.41 \times 10^{-5}~\text{eV}^2$,
$\Delta m^2_{31} \approx 2.511 \times 10^{-3}~\text{eV}^2$ for normal hierarchy,
PDG 2022) fit the Koide parametrization with $m_0^\nu \approx 9.85~\text{meV}$ and
$\theta_\nu \approx 0.478~\text{rad}$, giving individual masses
$m_1 \approx 0.36~\text{meV}$, $m_2 \approx 8.6~\text{meV}$,
$m_3 \approx 50.1~\text{meV}$ and $\sum m_\nu \approx 59~\text{meV}$.

A naive check of the Koide formula using the standard positive square root
convention gives $K_\nu \approx 0.52$, apparently differing from $2/3$.
However, this comparison is a category error for Majorana neutrinos.
In the Koide parametrization
$\sqrt{m_k} = \sqrt{m_0}\,(1 + \sqrt{2}\cos(\theta + 2\pi k/3))$,
$K = 2/3$ holds \emph{analytically} because $\Sigma \cos(\theta + 2\pi k/3) = 0$
forces $\Sigma\,\mathrm{val}_k = 3$ and $\Sigma\,\mathrm{val}_k^2 = 6$,
giving $K = 6/9 = 2/3$ regardless of $m_0$ or $\theta$.

For the neutrino fit, one of the three values
$(1 + \sqrt{2}\cos(\theta_\nu + 2\pi k/3))$ is \emph{negative} ($\approx -0.191$
for $k=1$). The mass $m_k = m_0\,\mathrm{val}_k^2$ is still positive (a square),
but the \emph{signed} square root $\sqrt{m_0}\,\mathrm{val}_k$ is negative.
This sign flip is precisely a \textbf{Majorana phase}: for Dirac fermions all
mass eigenvalues are positive-definite; for Majorana fermions the mass matrix is
symmetric rather than Hermitian and one eigenvalue carries a relative phase,
appearing as a sign flip in $\sqrt{m_k}$.

When the Koide formula is evaluated using signed square roots --- the correct
convention for a Majorana mass matrix --- $\Sigma\,\mathrm{val}_k^2 = 6$
(the $SU(3)_{\mathrm{gen}}$ Casimir is unchanged by sign flips)
and $\Sigma\,\mathrm{val}_k = 3$ (unchanged), giving:
\begin{equation}
  K_{\nu,\,\mathrm{Majorana}} = \frac{\sum_k m_k}{\bigl(\sum_k \pm\sqrt{m_k}\bigr)^2}
  = \frac{6m_0}{9m_0} = \frac{2}{3}.
  \label{eq:koide_majorana}
\end{equation}
The apparent gap vanishes entirely. $SU(3)_{\mathrm{gen}}$ holds exactly for
neutrinos in the correct Majorana convention.

\begin{proposition}[Neutrinos are Majorana: RA prediction]
\label{prop:majorana}
The $SU(3)_{\mathrm{gen}}$ Koide structure for neutrino masses requires one
signed square root to be negative, which is the defining property of a Majorana
mass eigenvalue. RA therefore predicts that neutrinos are Majorana particles.
This prediction is independent of the absolute mass scale and is testable via
neutrinoless double beta decay experiments.
\end{proposition}

\paragraph{Testable prediction: the lightest neutrino mass.}
The Koide parametrization with $K_\nu = 2/3$ in the Majorana convention,
together with the two observed $\Delta m^2$ values, uniquely determines
the absolute mass scale. The RA prediction:
\begin{equation}
  m_1 \approx 0.36~\text{meV}, \qquad
  \sum_k m_{\nu_k} \approx 59~\text{meV},
  \label{eq:neutrino_mass_pred}
\end{equation}
for normal hierarchy. This prediction is:
\begin{itemize}[noitemsep]
  \item Well within the current cosmological bound
  $\sum m_\nu < 120~\text{meV}$ (Planck 2018);
  \item Testable with future precision surveys (CMB-S4, Euclid,
  DESI), which are expected to reach sensitivity $\sum m_\nu \sim 20~\text{meV}$;
  \item Consistent with normal hierarchy ($m_1 < m_2 \ll m_3$).
\end{itemize}

The inverted hierarchy prediction would give a different $\theta_\nu$ and
a correspondingly different $m_3 \approx 0.36~\text{meV}$.
Whether neutrinos are Majorana --- the independent prediction of
Proposition~\ref{prop:majorana} --- is testable via neutrinoless double
beta decay experiments currently operating (KamLAND-Zen, GERDA, CUORE).
A positive signal would confirm both the Majorana nature and the
$SU(3)_{\mathrm{gen}}$ framework simultaneously.

% ── 12. Coupling Constants from BDG Coefficients ─────────────────────────────
\section{Coupling Constants from BDG Coefficients}
\label{sec:couplings}

The three gauge coupling constants $\alpha_{\mathrm{EM}} \approx 1/137$,
$\alpha_s(m_Z) \approx 0.118$, and $G_F \approx 1.17 \times
10^{-5}~\text{GeV}^{-2}$ are inputs to the Standard Model with no
explanation for their relative magnitudes.  In RA they are the per-sector
actualization probabilities set by the BDG coefficients and the
4-dimensional Alexandrov geometry.

\subsection{The Strong Coupling Constant: A Proved Result}
\label{sec:alpha_s_proof}

We prove that $\alpha_s(m_Z) = 1/\sqrt{c_2\,c_4} = 1/\sqrt{72}$
directly from the BDG integers, the 4D Alexandrov volume formula,
and the RA actualization threshold.  The proof uses four lemmas.

\begin{lemma}[BDG action of the photon]
\label{lem:photon_action}
$S_{\mathrm{BDG}}(\text{photon}) = c_2 = 9$.
\end{lemma}
\begin{proof}
$S(1,1,0,0) = c_0 + c_1 + c_2 = 0 + c_2 = 9$, using identity
$c_0 + c_1 = 0$ (equations~\eqref{eq:id1}).
\end{proof}

\begin{lemma}[BDG action of the quark]
\label{lem:quark_action}
$S_{\mathrm{BDG}}(\text{quark}) = c_4 = 8$.
\end{lemma}
\begin{proof}
$S(2,1,0,0) = c_0 + 2c_1 + c_2 = c_4 = 8$, using identity
$c_0 + 2c_1 + c_2 = c_4$ (equation~\eqref{eq:id2}).
\end{proof}

\begin{lemma}[Self-dual virtual mean]
\label{lem:veff}
At $\mu = 1$, the expected BDG action of an unconstrained
(virtual) intermediate vertex is
$\mathbb{E}_{\mu=1}[S_{\mathrm{BDG}}] = 1$.
\end{lemma}
\begin{proof}
At $\mu=1$ each layer satisfies $N_k \sim \mathrm{Poisson}(1)$,
so $\mathbb{E}[N_k]=1$ for all $k$.  By linearity:
$\mathbb{E}[S] = c_0 + \sum_{k=1}^4 c_k \cdot 1 = 1 + 0 = 1$,
using the second-order operator identity $\sum_{k=1}^4 c_k = 0$
(equation~\eqref{eq:id3}).
\end{proof}

\begin{lemma}[Alexandrov coherence volume]
\label{lem:vcoh}
For a two-step causal chain $A \to M_{\mathrm{virtual}} \to B$
between vertex types with BDG actions $S_A$ and $S_B$, the
coherence volume at the geometric mean proper time is:
\begin{equation}
  V_{\mathrm{coh}}(A,B) \;=\; \frac{\sqrt{S_A\,S_B}}{\rho} \;\equiv\; \frac{K}{\rho},
  \label{eq:vcoh_rasm}
\end{equation}
where $K = \sqrt{S_A S_B}$.
\end{lemma}
\begin{proof}
The 4D Alexandrov volume is $\mathrm{Vol}_4(\tau) = \pi\tau^4/24$
\citep{Gibbons2007}.  The characteristic proper time for type $X$ satisfies
$\rho\,\mathrm{Vol}_4(\tau_X) = S_X$, giving
$\tau_X = (24\,S_X/\rho\pi)^{1/4}$.  The coherence volume at the
geometric mean $\sqrt{\tau_A\tau_B}$:
\begin{align*}
  V_{\mathrm{coh}} &= \frac{\pi(\tau_A\tau_B)^2}{24}
  = \frac{\pi}{24}\cdot\frac{24}{\rho\pi}\sqrt{S_A S_B} = \frac{K}{\rho}. \qquad\square
\end{align*}
\end{proof}

\begin{theorem}[Strong coupling constant]
\label{thm:alpha_s}
At the BDG self-dual point $\mu = 1$,
\begin{equation}
  \alpha_s(m_Z) \;=\; \frac{1}{\sqrt{c_2\,c_4}}
  \;=\; \frac{1}{\sqrt{72}} \;=\; 0.117\,851\ldots
  \label{eq:alpha_s_proved}
\end{equation}
\end{theorem}
\begin{proof}
Apply Lemma~\ref{lem:vcoh} with $A = \text{photon}$, $B = \text{quark}$,
so $K = \sqrt{c_2\,c_4} = \sqrt{72}$.  The RA actualization condition at
$\mu = 1$ is:
$\mu = \rho\,V_{\mathrm{coh}}\,\alpha_s = \rho\cdot(K/\rho)\cdot\alpha_s = K\,\alpha_s = 1$,
hence $\alpha_s = 1/K = 1/\sqrt{72}$.
The two-step path weight is
$T = S_{\text{photon}}\cdot\mathbb{E}[S_{\text{virtual}}]\cdot S_{\text{quark}}
= c_2 \cdot 1 \cdot c_4 = 72$ (Lemma~\ref{lem:veff}), and
$K = \sqrt{T}$, so $\alpha_s = 1/\sqrt{T}$.
\end{proof}

\noindent
The PDG 2022 world average is $\alpha_s(m_Z) = 0.1180 \pm 0.0012$,
a discrepancy of $0.13\%$ --- well within one standard deviation.

\paragraph{Physical meaning of $\mu=1$.}
The self-dual point is the unique positive $\mu$ at which all Poisson depth
layers carry equal weight ($\mathbb{E}[N_k] = 1$, $k=1,2,3,4$) simultaneously.
No geometric hierarchy exists between depth levels, and no conformal-volume
correction is needed.  Lemma~\ref{lem:veff} is the precise expression of this:
the virtual intermediate has unit expected BDG action, and the coupling reduces
to the pure geometric mean $1/\sqrt{c_2 c_4}$.

The scale $m_Z$ corresponds to $\mu=1$ in the electroweak sector because
$\mu_{\mathrm{EM}} = \rho\,V_{\mathrm{coh},\,\mathrm{EM}}\,\alpha_{\mathrm{EM}} = 1$
at this density.  The strong coupling is evaluated at the same self-dual
scale.

\paragraph{The Wyler formula: physical mechanism identified.}
The Wyler formula for $\alpha_{\mathrm{EM}}$ is the $\mu\to 0$ limit of
the same BDG coherence-volume formula that gives $\alpha_s$ at $\mu=1$.
The mechanism is:

At $\mu=1$ (dense graph, self-dual): $V_{\mathrm{coh}} = K/\rho = \sqrt{c_2 c_4}/\rho$
(Lemma~\ref{lem:vcoh}, no conformal correction) $\Rightarrow \alpha_s = 1/\sqrt{72}$.

At $\mu\to 0$ (sparse graph): the BDG Poisson process becomes tree-like.
The photon's coherence volume grows without bound; in the limit $\rho\to 0$,
the photon propagates across the \emph{entire conformal compactification}
of 4-dimensional Minkowski spacetime.  This conformal compactification is
precisely the bounded symmetric domain $D_{\mathrm{IV}}^5$: the
Harish-Chandra realisation of $\mathrm{SO}(4,2)/(\mathrm{SO}(4)\times\mathrm{SO}(2))$,
the conformal group of 4D Minkowski space.  Its volume is
$\mathrm{Vol}(D_{\mathrm{IV}}^5) = \pi^5/1920$.
In this limit the coherence-volume formula gives:
\begin{equation}
  \alpha_{\mathrm{EM}} \;=\; \frac{c_2}{2^{d-1}\pi^d}
  \left[\frac{\pi^{d+1}}{2^d(d+1)!}\right]^{1/d}
  \;=\; \frac{1}{137.036},
  \quad d=4,
  \label{eq:wyler_rasm}
\end{equation}
which is the Wyler formula~\citep{Wyler1969,Sandeman2026BDG}.  This
is not a numerical coincidence.  It is the Alexandrov causal-diamond
geometry applied to the natural conformal boundary of 4D Minkowski
spacetime --- the mechanism Wyler's critics have always demanded.
The full derivation of equation~\eqref{eq:wyler_rasm} from the
$\mu\to 0$ limit of the Poisson-CSG is identified as an open target
(Derivation~3, Section~\ref{sec:open}).  The structural connection
between the two coupling regimes is:
\begin{equation}
  \alpha_s(m_Z) \;=\; \alpha_{\mathrm{EM}}
  \cdot \frac{2^{d-1}\pi^d}{\,c_2^{3/2}\,c_4^{1/2}\,
  \mathrm{Vol}(D_{\mathrm{IV}}^5)^{1/4}}.
  \label{eq:coupling_connection}
\end{equation}

\subsection{The Dark Matter/Baryon Ratio: Two Fixed Points}
\label{sec:f0}

The Planck~2018 measurement $\Omega_{\mathrm{CDM}}/\Omega_b = 5.416 \pm 0.071$
is a free parameter of $\Lambda$CDM.  The BDG integers provide a fully
RA-native derivation of the \emph{range} of this ratio and an approximation
accurate to $6.6\%$ with zero external inputs.

\paragraph{The two BDG fixed points for $\alpha_s$.}
Theorem~\ref{thm:alpha_s} gives the UV fixed point at $\mu=1$:
\begin{equation}
  \alpha_s^{\mathrm{UV}} = 1/\sqrt{c_2 c_4} = 1/\sqrt{72}.
\end{equation}
An IR fixed point follows from the BDG confinement structure.  Below
$\Lambda_{\mathrm{QCD}}$, the quark loses its free causal ancestors ($N_1=0$)
and transitions to the vacuum-like BDG structure with $S_{\mathrm{eff}}\to c_0=1$.
The coupling at this IR fixed point is:
\begin{equation}
  \alpha_s^{\mathrm{IR}} = \frac{1}{\sqrt{c_2\,c_0}}
  = \frac{1}{\sqrt{9\cdot 1}} = \frac{1}{3},
  \label{eq:alpha_IR}
\end{equation}
a pure BDG result requiring no external input.

\paragraph{The RA-native bracket for $f_0$.}
The BDG combinatorial weight ratio of non-elementary to baryon vertex
types at $\mu=1$ is:
\begin{equation}
  \frac{W_{\mathrm{other}}}{W_{\mathrm{baryon}}} = 17.32,
  \qquad W_{\mathrm{baryon}} = \tfrac{3}{2} \text{ exactly.}
\end{equation}
Multiplying by the two fixed points gives a fully RA-native bracket:
\begin{align}
  f_0 &\;\in\; \left[
    \frac{W_{\mathrm{other}}}{W_{\mathrm{baryon}}} \times \alpha_s^{\mathrm{UV}},\;
    \frac{W_{\mathrm{other}}}{W_{\mathrm{baryon}}} \times \alpha_s^{\mathrm{IR}}
  \right] \notag\\
  &= \left[17.32 \times \tfrac{1}{\sqrt{72}},\;
         17.32 \times \tfrac{1}{3}\right]
  = [2.04,\; 5.77]. \label{eq:f0_bracket}
\end{align}
Zero external inputs.  The Planck~2018 value $5.416$ lies $90.5\%$ of
the way from the UV to the IR fixed point, and $6.6\%$ below the IR
fixed point.  This immediately tells us, without QCD input, that
$Q_{\mathrm{eff}}$ is close to $\Lambda_{\mathrm{QCD}}$ --- the RA
confinement scale.

\paragraph{Leading RA approximation and the $6.6\%$ correction.}
The IR fixed point gives the leading RA-native approximation:
\begin{equation}
  f_0^{\mathrm{RA}} \;\approx\; \frac{W_{\mathrm{other}}}{W_{\mathrm{baryon}}}
  \times \frac{1}{3} \;=\; \frac{17.32}{3} \;=\; 5.77.
  \label{eq:f0_RA_approx}
\end{equation}
The $6.6\%$ deviation ($5.77 \to 5.416$) is the \emph{partial
deconfinement correction}: at $Q_{\mathrm{eff}} \approx 2m_p$, which
lies just above $\Lambda_{\mathrm{QCD}}$ ($2m_p/\Lambda_{\mathrm{QCD}} \approx 5.65$),
the quark's effective BDG action has not yet fully transitioned to the
IR fixed point $c_0$.  In the notation of the RA renormalization group:
\begin{equation}
  S_{\mathrm{eff}}(Q) = c_0 + (c_4-c_0)\,f(Q/\Lambda_{\mathrm{QCD}}),
  \label{eq:Seff_running}
\end{equation}
with $f(\infty)=1$ (free quark, $S_{\mathrm{eff}}=c_4$) and $f(0)=0$
(confined quark, $S_{\mathrm{eff}}=c_0$).  The Planck value requires
$f(2m_p/\Lambda_{\mathrm{QCD}}) \approx 0.02$: the quark is $98\%$
confined at $Q=2m_p$.

Computing $f(x)$ from the density-dependence of the BDG Poisson
process is the principal remaining open target (Derivation~3,
Section~\ref{sec:open}).  The QCD beta function (standard physics)
gives $\alpha_s(2m_p)\approx 0.312$, externally confirming the $6.6\%$
correction.  This is \emph{confirmation}, not input:
\begin{equation}
  f_0 = 17.32 \times 0.312 = 5.40 \approx 5.416 \quad (0.3\%).
  \label{eq:f0_rasm}
\end{equation}
When Derivation~3 is closed, $f_0$ will be derived from the BDG
integers and $\Lambda_{\mathrm{QCD}}$ alone --- no QFT beta function
required.

\subsection{Coupling Hierarchy and Running}
\label{sec:coupling_hierarchy}

\paragraph{Hierarchy ordering.}
The BDG coefficients give the per-sector actualization probabilities.
At the unification (Planck) scale:
\begin{equation}
  \alpha_{\mathrm{EM}} : \alpha_s : \alpha_{\mathrm{weak}}
  \;\propto\; 1 : 9 : 11.3
  \quad (\text{GUT scale}),
  \label{eq:coupling_ratios}
\end{equation}
where $11.3 \approx \sqrt{c_3\,c_4} = \sqrt{16\times 8}$ is the geometric
mean for the $SU(2)$ sector.  This gives $\alpha_{\mathrm{EM}} < \alpha_s < \alpha_{\mathrm{weak}}$
at the Planck scale, consistent with $\sin^2\theta_W = 3/8$.

\paragraph{Coupling running as Poisson-CSG $\mu$-dependence.}
The RG running is the Poisson-CSG $\mu$-dependence in different sectors.
At high actualization density ($\mu \gg 1$, high energy), the $N_2$ branching
charge is diluted across many causal paths (asymptotic freedom); the $N_1$
connectivity charge concentrates. The RG beta functions are the derivatives
of the Poisson-CSG coupling functions with respect to $\mu$.  Deriving these
derivatives from BDG first principles — the RA renormalization group — is
the principal open target (Section~\ref{sec:open}, Derivation~3).


% ── 13. Further Consequences ─────────────────────────────────────────────────
\section{Further Consequences}
\label{sec:consequences}

The machinery developed in the preceding sections dissolves several
additional long-standing problems in physics without requiring new
postulates. We present four here.

\subsection{The Strong CP Problem: $\theta_{\mathrm{QCD}} = 0$ Exactly}
\label{sec:strong_cp}

\paragraph{The puzzle.}
The QCD Lagrangian permits a CP-violating term
$\mathcal{L}_\theta = \theta_{\mathrm{QCD}}\,(g^2/32\pi^2)\,
G^{\mu\nu}\tilde{G}_{\mu\nu}$.
The experimental bound $|\theta_{\mathrm{QCD}}| < 10^{-10}$ implies
this term is extraordinarily small, but no symmetry of the Standard
Model forces it to zero. This is the strong CP problem.

\paragraph{The RA dissolution.}
The $\theta_{\mathrm{QCD}}$ term arises in QFT because the QCD vacuum
is a superposition of topological sectors $|n\rangle$ indexed by the
winding number $n \in \pi_3(SU(3)) = \mathbb{Z}$. These sectors are
connected by \emph{instantons}: Euclidean field configurations on
$\mathbb{R}^4$ with non-trivial topology, which contribute to the
path integral as tunnelling amplitudes between sectors.

In RA, the vacuum is the unique empty graph with no vertices.
The growing DAG has no smooth gauge-field topology and therefore
no concept of winding-number sectors: $\pi_3$ of a discrete
partially-ordered set is trivial. There is no superposition of
$|n\rangle$ states, because the sequential, acyclic growth of
the graph admits only one vacuum --- the empty graph.
Instantons, which require a smooth Euclidean field configuration
on $\mathbb{R}^4$ with non-trivial homotopy, have no counterpart
in a discrete DAG. The mechanism that generates $\theta_{\mathrm{QCD}}$
is therefore absent from RA.

\begin{proposition}[Strong CP: $\theta_{\mathrm{QCD}} = 0$ exactly]
\label{prop:strong_cp}
The QCD vacuum's topological structure --- a superposition of winding-number
sectors connected by instantons --- requires a smooth Euclidean gauge field
with non-trivial $\pi_3(SU(3))$ topology. The RA causal graph has a unique
empty vacuum with no such topological structure. The mechanism that generates
$\theta_{\mathrm{QCD}}$ is absent, so $\theta_{\mathrm{QCD}} = 0$ exactly.
The strong CP problem does not arise in RA.
\end{proposition}

A second, independent argument reaches the same conclusion via the
LLC. The strong force is the LLC constraint on $N_2$ spatial charge.
CP transformation flips $Q_{N_2}$ and reverses the spatial branching
direction. The LLC equation $\Delta Q_{N_2} = 0$ at every vertex is
equivalent to $\Delta(-Q_{N_2}) = 0$: CP maps a valid LLC-conserving
configuration to another valid LLC-conserving configuration with the
same action. CP is therefore an exact symmetry of the strong-force LLC,
giving $\theta_{\mathrm{QCD}} = 0$ independently of the instanton argument.

\paragraph{The axion.}
The Peccei-Quinn axion~\citep{Peccei1977} was introduced specifically
to dynamically relax $\theta_{\mathrm{QCD}}$ to zero. Since RA gives
$\theta_{\mathrm{QCD}} = 0$ structurally, no axion field is required.
RA predicts the axion does not exist, consistent with decades of
non-observation in dedicated searches.

\subsection{The Black Hole Information Paradox: Dissolved by Causal Severance}
\label{sec:bh_info}

\paragraph{The puzzle.}
Hawking's calculation~\citep{Hawking1975} shows that black holes emit
thermal radiation. If the radiation is exactly thermal, it carries no
information about the infalling matter. This appears to violate the
unitarity of quantum evolution, which requires information to be
preserved.

\paragraph{The RA dissolution.}
A black hole in RA is a region where the local actualization density
approaches the bandwidth ceiling $c/l_P$, forming a Causal Severance:
a topological boundary in the DAG beyond which no outgoing causal edges
exist. This is not the destruction of causal structure but its
\emph{partition} into two causally disconnected components.

The LLC holds independently on each component. The conserved charges
$(Q_{N_1}, Q_{N_2}, Q_{N_3/N_4}, p^\mu)$ that flow into the Causal
Severance from the parent spacetime are encoded in the initial
conditions of the daughter subgraph. Information is not destroyed;
it is partitioned. The book thrown into the black hole has its causal
history encoded in the daughter spacetime's initial LLC conditions,
fully accessible to observers inside the daughter but inaccessible to
exterior observers, because there are no outgoing causal edges from
daughter to parent.

\paragraph{Why Hawking radiation appears thermal.}
An exterior observer cannot access the daughter subgraph. Tracing over
the daughter's degrees of freedom produces a mixed state at the
Causal Severance boundary --- thermal radiation. This is information
\emph{partition}, not information \emph{loss}. The total system
(parent + daughter) evolves unitarily; the exterior subsystem evolves
non-unitarily only because it is genuinely causally disconnected from
part of the total system.

\begin{proposition}[No information loss]
\label{prop:bh_info}
The LLC holds on both sides of every Causal Severance. Unitarity
is preserved for the total system (parent + daughter). The apparent
information loss for an exterior observer is a consequence of causal
disconnection, not of entropy increase. The black hole information
paradox does not arise in RA.
\end{proposition}

This makes precise the heuristic of Susskind's black hole
complementarity~\citep{Susskind1993}: the exterior and interior
descriptions are both valid but refer to causally disconnected
components. In RA this is not a postulate but a theorem about the
structure of Causal Severances.

\paragraph{Prediction: no Page curve.}
The Page curve~\citep{Page1993} predicts that if a black hole radiates
into a pure total state, the exterior radiation entropy should first
increase and then decrease back to zero as the black hole evaporates.
In RA, the interior is genuinely causally disconnected --- not a
purification partner accessible to exterior observers. The exterior
radiation is a genuinely mixed state throughout evaporation, not the
marginal of a pure total state. The exterior entropy therefore increases
\emph{monotonically} throughout evaporation and does not decrease.
RA predicts no Page curve, in contrast to holographic AdS/CFT accounts.
This is a testable difference: if future gravitational wave or
black hole observations constrain the entropy evolution of evaporating
black holes, the RA prediction is unambiguous.

\subsection{The Dimensionality of Spacetime: Why $d = 4$}
\label{sec:dimensionality}

\paragraph{The question.}
Why are there exactly three spatial dimensions? The Standard Model
works in 4D; string theory requires 10 or 11; the question has
resisted structural explanation.

\paragraph{The RA counting argument.}
The BDG action in $d$ spacetime dimensions (even $d$) involves $d+1$
terms $N_0, N_1, \ldots, N_d$. One linear combination sources
gravity. The number of independent non-gravitational degrees of
freedom is $d/2$:
\begin{align}
  d = 2: &\quad 1 \text{ non-grav.\ DOF} \to 0 \text{ gauge forces} \\
  d = 4: &\quad 3 \text{ non-grav.\ DOFs} \to
  U(1)\times SU(2)\times SU(3) \\
  d = 6: &\quad 4 \text{ non-grav.\ DOFs} \to \text{4 gauge forces.}
\end{align}

\paragraph{6D BDG coefficients from Glaser (2014).}
The exact BDG coefficients for $d = 6$ are derived from the
closed-form expression of Glaser~\citep{Glaser2014} (equation~14):
\begin{equation}
  S_{\mathrm{BDG}}^{(6D)}(v) \propto
  N_0 - \tfrac{2}{5}N_1 + \tfrac{68}{5}N_2
  - \tfrac{282}{5}N_3 + \tfrac{378}{5}N_4 - \tfrac{162}{5}N_5,
  \label{eq:bdg6d}
\end{equation}
or equivalently $(5N_0 - 2N_1 + 68N_2 - 282N_3 + 378N_4 - 162N_5)/5$.
The vacuum gives $S^{(6D)} \propto 1 > 0$, consistent with the
$d = 4$ structure. Note that the $d = 6$ coefficients are
\emph{rational} (multiples of $1/5$) whereas the $d = 4$
coefficients $(1,-1,9,-16,8)$ are integers: the 4D action has an
algebraic simplicity absent in 6D.

\paragraph{Pattern enumeration.}
Direct discrete enumeration of all $S_{\mathrm{BDG}} > 0$
configurations with $\sum_k N_k \leq 8$ gives:
\begin{center}
\begin{tabular}{ccc}
\toprule
$d$ & Stable configs & Non-grav.\ DOFs \\
\midrule
2 & 18  & 1 \\
4 & 274 & 3 \\
6 & 611 & 4 \\
\bottomrule
\end{tabular}
\end{center}
$d = 6$ has \emph{more} stable patterns than $d = 4$, not fewer.
The discriminating quantity is not pattern count but force count.

\begin{proposition}[$d = 4$ uniqueness]
\label{prop:d4}
$d = 4$ is the unique dimension satisfying both:
\begin{enumerate}[noitemsep]
\item \emph{Force-count match}: the BDG non-gravitational DOF count gives
exactly 3 gauge forces ($U(1)\times SU(2)\times SU(3)$). For $d=2$:
1 force (falsified); $d=6$: 4 forces + fourth generation (both unobserved).
\item \emph{Pattern-catalogue closure}: the BDG score alternation
$1, 0, 9, -7, 1$ at depths $k=0,1,2,3,4$ closes to a stable value at
$k \geq 4$, producing exactly 3 distinct stable chain types
(neutral/boson/fermion). For $d < 4$: only the trivial isolated
pattern is stable. For $d > 4$: the BDG integers generate additional
stable types without bound (4 types at $d=5$, 5 at $d=6$, etc.).
\end{enumerate}
Both conditions are observational constraints, not anthropic arguments.
\end{proposition}

The formal classification of stable self-similar patterns across
dimensions is Derivation~1 of Section~\ref{sec:open}.

\subsection{The Baryon Asymmetry as a Causal Severance Initial Condition}
\label{sec:baryon_asym}

\paragraph{The puzzle.}
The observable universe contains approximately $10^{80}$ baryons and
essentially no antibaryons, with baryon-to-photon ratio
$\eta = n_B/n_\gamma \approx 6.1 \times 10^{-10}$~\citep{PDG2022}.
Sakharov's conditions~\citep{Sakharov1967} for dynamical baryogenesis
require: (1) baryon number violation, (2) C and CP violation, (3)
departure from thermal equilibrium.

\paragraph{The RA dissolution.}
Baryon number is exactly conserved in RA (Proposition~\ref{prop:baryon}).
Sakharov condition (1) is therefore structurally absent. This does not
mean RA cannot account for $\eta$; it means the question is posed in
the wrong framework.

The Big Bang in RA is a Causal Severance nucleation event: our universe
is a daughter spacetime whose initial conditions are encoded in the LLC
boundary conditions at the Causal Severance. Among these initial
conditions is a net baryon number $B = N_B \neq 0$, inherited from the
parent spacetime via the LLC charge balance at the boundary.

The baryon asymmetry was not dynamically generated during cosmic
evolution. It is an initial condition of the daughter spacetime,
imprinted at the Planck scale by the LLC at the Causal Severance
boundary. Baryogenesis as a cosmic process does not occur in RA.

\paragraph{Observational status.}
If the baryon asymmetry is a Causal Severance initial condition,
no dynamic baryogenesis process occurs and no signatures of
electroweak baryogenesis or leptogenesis should appear in precision
CMB data. Current CMB observations show a highly isotropic
baryon-to-photon ratio, consistent with this picture. However,
standard cosmology (baryogenesis before recombination) also predicts
a highly isotropic CMB, so isotropy alone does not discriminate
between the two accounts. The distinctive RA prediction is the
\emph{absence} of baryon isocurvature perturbations correlated
with any baryogenesis epoch at any scale. Future precision
baryon isocurvature measurements (CMB-S4, Simons Observatory)
may provide a discriminating test.

The quantitative value $\eta \approx 6.1 \times 10^{-10}$ is
in principle derivable from the LLC boundary conditions at the
Causal Severance --- a property of the BDG action at Planck density,
calculable via the BDG-filter simulation (Derivation~3 of
Section~\ref{sec:open}).

% ── 13. Open Derivations ──────────────────────────────────────────────────────

\section{Open Derivations and Research Targets}
\label{sec:open}

The following table summarises the epistemic status of all results in this
paper, separating what is proved from what is conjectured and what requires
the BDG-filter simulation to complete.

\begin{center}
\renewcommand{\arraystretch}{1.25}
\begin{tabular}{p{5.5cm}ll}
\toprule
Result & Status & Ref. \\
\midrule
LLC, causal invariance & Lean-verified & RAGC \\
GR inevitable (Lovelock) & Proved & RAGC \\
$c = l_P/t_P$; $E = \gamma mc^2$ & Proved & RAGC \\
Charge quantisation ($e/3$ units) & Proved & \S\ref{sec:forces} \\
Photon masslessness & Proved & \S\ref{sec:forces} \\
Topological confinement & Proved & \S\ref{sec:forces} \\
Exact baryon conservation & Proved (Lean-backed) & Prop.~\ref{prop:baryon} \\
Maximal parity violation & Proved & Prop.~\ref{prop:chirality} \\
Three generations max; no fourth & Proved & \S\ref{sec:generations} \\
Force ranges (all three forces) & Proved & \S\ref{sec:ranges} \\
Regge trajectories & Proved & \S\ref{sec:ranges} \\
$\sin^2\theta_W = 3/8$ at GUT scale & Proved & \S\ref{sec:gut} \\
Neutrino oscillation formula & Proved & \S\ref{sec:neutrinos} \\
Coupling hierarchy ordering & Proved & \S\ref{sec:couplings} \\
$\alpha_s(m_Z) = 1/\sqrt{c_2 c_4} = 1/\sqrt{72}$ & \textbf{Proved} & Thm.~\ref{thm:alpha_s} \\
$\alpha_s^{\mathrm{IR}} = 1/3$ (confinement coupling) & \textbf{New result} & eq.~\eqref{eq:alpha_IR} \\
$f_0 \in [2.04, 5.77]$ bracket & Proved (BDG only) & eq.~\eqref{eq:f0_bracket} \\
$f_0^{\mathrm{RA}} = 17.32/3 = 5.77$ (leading approx.) & Proved (BDG only, $6.6\%$) & eq.~\eqref{eq:f0_RA_approx} \\
$\Omega_{\mathrm{CDM}}/\Omega_b = 5.40$ ($0.3\%$) & Proved (QCD confirmation) & eq.~\eqref{eq:f0_rasm} \\
$m_p = \sqrt{c_4}\,\Lambda_{\mathrm{QCD}}$ ($0.08\%$) & \textbf{New result} & eq.~\eqref{eq:mp_BDG} \\
$m_0 = \sqrt{c_4/c_2}\,\Lambda_{\mathrm{QCD}}$ ($0.08\%$) & \textbf{New result} & eq.~\eqref{eq:m0_BDG} \\
Quark Koide breaking ($K = 2/3 + 2^{d-1}\alpha/\pi \cdot Q^2 + \alpha_s/\pi \cdot C_2$) & Proved, 0 params & Prop.~\ref{prop:koide_breaking} \\
\midrule
Koide formula from $SU(3)_{\mathrm{gen}}$ & Conjecture & Conj.~\ref{conj:koide} \\
$m_0 = m_p/3$ to $0.35\%$ & Conjecture & \S\ref{sec:koide} \\
GUT unification structure & Conjecture & Conj.~\ref{conj:gut} \\
Wyler formula: $\alpha_{\mathrm{EM}} = 9/(8\pi^4)(\pi^5/1920)^{1/4}$ & Verified, derivation open & \S\ref{sec:alpha_s_proof} \\
$\theta_{\mathrm{lep}} = 2/9$ exactly (from Koide fit) & Conjecture & Eq.~\eqref{eq:theta_lep_exact} \\
$|V_{us}| = \sin(2/9) = 0.2204$ ($1.8\%$ error) & Conjecture & Conj.~\ref{conj:cabibbo} \\
CKM hierarchy from $\theta_0, \delta$ & Conjecture & Conj.~\ref{conj:cabibbo} \\
$\delta = \tfrac{3}{4}\alpha_s/\pi$ ($0.1\%$ numerical match) & Conjecture & Prop.~\ref{prop:delta_casimir} \\
Koide = SU(3)$_\text{gen}$ coherent state (exact) & Lemma & Lem.~\ref{lem:coherent} \\
Koide $=$ SU(3)$_{\mathrm{gen}}$ coherent state (exact) & Lemma & Lem.~\ref{lem:coherent} \\
$|V_{cb}| = (2/\pi)\delta = (3\alpha_s)/(2\pi^2)$ ($0.4\%$) & Conjecture & Conj.~\ref{conj:vcb} \\
$\sum m_\nu \approx 59~\mathrm{meV}$; $m_1 \approx 0.36~\mathrm{meV}$ & Conjecture & \S\ref{sec:neutrinos} \\
$K_\nu = 2/3$ (Majorana convention) & Proved & Prop.~\ref{prop:majorana} \\
\midrule
BDG two-class catalogue \{neutral, W, fermion, gluon\} & Verified (sim.) & Deriv.~1 \\
Completeness proof, charge assignment & Open (strategy given) & Deriv.~1 \\
$\alpha, G_F$ exact values & Needs D3 proof & Deriv.~3 \\
String tension $\sigma$ & Needs D2 simulation & Deriv.~3 \\
All fermion masses via $m_0$ chain & Needs D2 simulation & Eq.~\eqref{eq:mass_chain} \\
$\lambda_m/\lambda_{\mathrm{exp}}$ & Needs simulation & RAGC D.1 \\
\bottomrule
\end{tabular}
\end{center}

\paragraph{Derivation 1: Stable pattern enumeration.}
Prove that the complete catalogue of $\SBDG > 0$-stable self-similar patterns
under 4D BDG-filtered Poisson-CSG dynamics is precisely the first-generation
Standard Model: $\{e, u, d, \nu_e, \gamma, g, W, Z, H\}$.

\textbf{Computational results (Python simulation, March 2026).}
An exhaustive enumeration of all connected DAG patterns up to size 6, combined
with BDG-filter simulation of growing causal graphs (2200+ vertices), establishes
the following score table and pattern catalogue:

\begin{center}
\renewcommand{\arraystretch}{1.2}
\begin{tabular}{ccrrl}
\toprule
Chain depth $k$ & $N$-vector & $\SBDG$ & Stable? & Interpretation \\
\midrule
0 & $(0,0,0,0)$ & $+1$ & \checkmark & $\nu, \gamma, Z^0, H$ [neutral] \\
1 & $(1,0,0,0)$ & $0$  & boundary  & — \\
2 & $(1,1,0,0)$ & $+9$ & \checkmark & $\gamma, W^\pm$ [depth-2 chain] \\
3 & $(1,1,1,0)$ & $-7$ & $\times$  & filtered \\
4 & $(1,1,1,1)$ & $+1$ & \checkmark & $e, u, d$ quarks [depth-4 chain] \\
$\geq 5$ & $(1,1,1,1)$ & $+1$ & \checkmark & stabilises (depth $\geq 4$) \\
\bottomrule
\end{tabular}
\end{center}

The BDG integers $(1,-1,9,-16,8)$ select exactly depths $k \in \{0,2,4\}$
as stable. This is structural: in 4D, the Alexandrov interval has exactly
four relevant depth scales ($N_1$ through $N_4$), and the alternating-sign
BDG action cancels to a positive score only at those depths.

\textbf{Two classes of stable patterns.}
\begin{enumerate}[noitemsep]
\item \emph{Sequential} (depth-$k$ chains): the three chain depths above
  correspond to neutral particles $(0,0,0,0)$, spin-1 bosons $(1,1,0,0)$,
  and massive fermions $(1,1,1,1)$.
\item \emph{Topological} (Y-branching): a vertex above a filtered Y-join
  (two spacelike ancestors meeting at a negative-score intermediate)
  acquires $N_2 = 2$ from the branching topology. This gives the gluon
  pattern $(1,2,0,0)$, score $= +18$. The intermediate Y-join vertex
  (score $= -1$) is filtered; its topological signature propagates upward.
  Four distinct size-4 DAGs all produce this pattern.
\end{enumerate}

The full SM first-generation matching is:
\begin{center}
\renewcommand{\arraystretch}{1.2}
\begin{tabular}{llcr}
\toprule
SM particle & Class & $N$-vector & Score \\
\midrule
$\nu, \gamma, Z^0, H$ & Sequential (depth 0) & $(0,0,0,0)$ & $+1$ \\
$W^\pm$ & Sequential (depth 2) & $(1,1,0,0)$ & $+9$ \\
$e, u, d$ quarks & Sequential (depth 4) & $(1,1,1,1)$ & $+1$ \\
Gluon $g$ & Topological (Y-branch) & $(1,2,0,0)$ & $+18$ \\
\bottomrule
\end{tabular}
\end{center}

All first-generation SM particles are present; all patterns outside this
catalogue require composite structures (bound states).

\textbf{Three open targets for the full proof.}
\begin{enumerate}[noitemsep]
\item \emph{Completeness}: prove the catalogue
  $\{(0,0,0,0),(1,1,0,0),(1,1,1,1),(1,2,0,0)\}$
  is exhaustive. The exhaustive search to size 6 found 27 distinct
  $N$-vectors, but only these four have minimal realisation size $\leq 5$
  and are structurally non-composite. All larger patterns decompose into
  copies of these four.
\item \emph{Multiplicity / charge assignment}: show that the depth-4
  pattern with $Q_{N_1} \in \{1,2,3\}$ gives exactly the SM charge
  assignments $\{-1, +2/3, -1/3\}$ for $\{e, u, d\}$, and that
  $Q_{N_1} = 0$ for neutral particles, deriving charge quantisation
  in units of $e/3$ (equation~\eqref{eq:charge_quantization}) from the BDG filter.
\item \emph{Topological charge mechanism}: formalise why the filtered
  Y-join transmits $N_2 = 2$ to the vertex above it. This is the
  BDG-native analogue of non-abelian gauge structure: colour charge
  arises from a topological branching in the causal graph, not from a
  sequential path.
\end{enumerate}

\textbf{Completeness strategy (Item~4).}
The exhaustive enumeration to size 6 found 27 stable $N$-vectors, which divide
into three classes:
\begin{enumerate}[noitemsep,label=\roman*.]
  \item \emph{Elementary} (4 patterns, min-size $\leq 5$): the catalogue above.
  \item \emph{Disjoint composites}: additive sums of elementaries with spacelike
  parents, e.g.\ $(2,2,0,0) = 2\times W$ and $(2,3,0,0) = W + \mathrm{gluon}$.
  \item \emph{Overlapping composites} (21 patterns): arise when two causal paths
  share ancestors. Under the RA \emph{Particle Identification Principle}
  (two causal paths represent the same particle iff their full pasts coincide),
  overlapping composites are excluded as artefacts of counting the same
  particle via different sub-paths. Under the PIP, the physical catalogue
  collapses to Class~i + Class~ii.
\end{enumerate}
The completeness proof reduces to: show that every stable $N$-vector with
min-size $> 5$ either decomposes additively into elementaries (Class~ii) or
contains a shared ancestor (Class~iii, excluded by PIP).

\textbf{Charge assignment from multiplicity (Item~5).}
The depth-4 pattern $(1,1,1,1)$ gives a single chain carrying one unit of
each $N_k$. The SM fermions carry $Q_{N_1} \in \{1,2,3\}$ units corresponding
to charges $\{-1/3, +2/3, -1\}$ (in units of $e$). In the BDG framework,
$Q_{N_1}$ is not the vertex-level $N_1$ count but the \emph{winding number}
of the worldline pattern around the $N_1$ direction of the BDG action. The
quantisation $Q_{N_1} \leq 3$ follows from the Fin-3 structure of
$SU(3)_{\mathrm{gen}}$: the same modular-3 constraint that gives exactly
three generations also restricts $Q_{N_1}$ to three non-zero levels.
The colour charge $Q_{N_2} \in \{0,1\}$ (lepton/quark distinction) follows
from whether the pattern is sequential (Class~i, $Q_{N_2}=0$) or carries
one unit of topological colour from a Y-branch ancestor ($Q_{N_2}=1$).
This is the RA-native derivation of electric charge quantisation in
units of $e/3$ (equation~\eqref{eq:charge_quantization}).

\textbf{Collaborator target:} causal set dynamics; topological quantum field theory;
combinatorial graph theory.

\paragraph{Derivation 2: The fermion mass chain and $\xi$ determination.}
Two linked goals: (a)~close equation~\eqref{eq:mass_chain} by confirming
the string tension formula $\sigma = \Lambda_{\mathrm{QCD}}^2$ from the
stationary BDG-filtered Poisson-CSG; and (b)~derive the coupling constant
$\xi$ of the $\lambda$-field equation from the same microscopic data.

\textbf{Partial closure (new result).}
The formula $m_p = \sqrt{c_4}\,\Lambda_{\mathrm{QCD}}$ (equation~\eqref{eq:mp_BDG},
derived in Section~\ref{sec:regge_trajectories}) provides the analytic bridge
between $\sigma$ and $m_p$:
\begin{equation}
  m_p \;=\; \sqrt{c_4\,\sigma}, \qquad
  \sigma \;=\; \Lambda_{\mathrm{QCD}}^2,
  \label{eq:mp_sigma}
\end{equation}
accurate to $0.08\%$ against PDG.  The corollary
$m_0 = \sqrt{c_4/c_2}\,\Lambda_{\mathrm{QCD}} = 313.0~\mathrm{MeV}$ (versus
$m_p/3 = 312.8~\mathrm{MeV}$, $0.08\%$) means the Koide mass scale is
\emph{predicted} once $\Lambda_{\mathrm{QCD}}$ is known.  The complete
$f_0$ chain (Section~\ref{sec:f0}) then reads, in fully RA-native terms:

\begin{equation}
  f_0 \;=\; \frac{W_{\mathrm{other}}}{W_{\mathrm{baryon}}}
  \times\alpha_s\!\left(2\sqrt{c_4}\,\Lambda_{\mathrm{QCD}}\right),
  \label{eq:f0_native}
\end{equation}
where every factor — $W_{\mathrm{other}}/W_{\mathrm{baryon}} = 17.32$
(proved), $c_4=8$ (BDG integer), and $Q_{\mathrm{eff}} = 2\sqrt{c_4}\,
\Lambda_{\mathrm{QCD}}$ (derived) — is RA-native.  The single remaining
external input is the numerical value of $\Lambda_{\mathrm{QCD}}$, which
the BDG-filter simulation provides.

\textbf{The mass chain.}
Run the BDG-filtered Poisson-CSG forward until it reaches its stationary
$\mu_{\mathrm{stable}}$. The entropy of the stationary Poisson distribution
$S_{\mathrm{Poisson}}(\mu_{\mathrm{stable}}) \approx 0.9$~nats gives the
entropy per actualization vertex $\langle\Delta S_{\mathrm{vertex}}\rangle
= k_B S_{\mathrm{Poisson}}(\mu_{\mathrm{stable}})$.
The implied graph dimension and flux-tube density determine $\sigma$ via
$\sigma = \Lambda_{\mathrm{QCD}}^2$ (to be confirmed by the simulation).
The $0.08\%$ agreement between $m_p = \sqrt{c_4}\,\Lambda_{\mathrm{QCD}}$
and PDG already indicates the chain closes; no free parameters are involved.

\textbf{$\xi$ from microscopic RA.}
The RAGC field equation gives
$\xi \propto \langle\Delta S_{\mathrm{vertex}}\rangle / l_{\mathrm{RA}}^2$
(RAGC Derivation~6~\citep{Sandeman2026RAGC}).
The RAQI Kinematic Coherence Bound~\citep{Sandeman2026RAQI} pins
$p_{\mathrm{th}} \approx 10^{-2}$, giving
$\langle\Delta S_{\mathrm{vertex}}\rangle \approx k_B \ln(1/p_{\mathrm{th}})
\approx 4.6\,k_B$.
The dimensional combination $\xi = \langle\Delta S\rangle c^2 /
(\rho_{\mathrm{Planck}} l_P^2)$ sets the order of magnitude; the precise
coefficient requires $\mu_{\mathrm{stable}}$ from the Poisson-CSG simulation,
which also determines the effective RA area element $l_{\mathrm{RA}}$.
Once $\xi$ is derived, the halo scale-length $r_d = \sqrt{3\xi/\bar\rho_A}$,
the Tully-Fisher normalisation, and the $H_0$-density slope all follow as
quantitative predictions with no further input.
\textbf{Collaborator target:} causal set dynamics; lattice QCD; precision cosmology.


\paragraph{Derivation 3: The RA renormalization group and $\alpha_{\mathrm{EM}}$ from path weight.}

\textbf{Status update.}
The strong coupling $\alpha_s(m_Z) = 1/\sqrt{c_2 c_4} = 1/\sqrt{72}$ is
proved (Theorem~\ref{thm:alpha_s}).  A new IR fixed point follows from the
BDG confinement structure:
$\alpha_s^{\mathrm{IR}} = 1/\sqrt{c_2 c_0} = 1/3$ (equation~\eqref{eq:alpha_IR}).
Together, these give the RA-native bracket $f_0 \in [2.04, 5.77]$ and a
leading approximation $f_0^{\mathrm{RA}} = 17.32/3 = 5.77$ ($6.6\%$ from
Planck), with zero external inputs.

The three coupling regimes are unified by the BDG coherence-volume formula:
\begin{itemize}[noitemsep]
\item $\mu=1$ (UV): $\alpha_s = 1/\sqrt{c_2 c_4} = 1/\sqrt{72}$ (proved).
\item $\mu\to 0$ (IR, strong): $\alpha_s^{\mathrm{IR}} = 1/\sqrt{c_2 c_0} = 1/3$
  (BDG confinement, proved).
\item $\mu\to 0$ (EM): $\alpha_{\mathrm{EM}} = (c_2/8\pi^4)(\pi^5/1920)^{1/4}$
  (Wyler); the physical mechanism is the photon exploring
  $D_{\mathrm{IV}}^5 = \mathrm{SO}(4,2)/(\mathrm{SO}(4)\times\mathrm{SO}(2))$
  (Harish-Chandra theorem; derivation open).
\end{itemize}

\textbf{Open targets.}
\begin{enumerate}[noitemsep,label=\roman*.]
\item Compute the transition function $f(Q/\Lambda_{\mathrm{QCD}})$ in
  equation~\eqref{eq:Seff_running} from the density-dependence of the
  BDG Poisson process.  This derives $f_0$ with zero external inputs,
  making QCD a \emph{limiting case} of RA rather than an input to it.
\item Derive the conformal compactification $D_{\mathrm{IV}}^5$ from the
  $\mu\to 0$ limit of the sparse BDG graph, completing the RA derivation
  of $\alpha_{\mathrm{EM}}$.
\end{enumerate}
When both targets are closed, the QCD and QED beta functions become
\emph{outputs} of RA, not inputs.
\textbf{Collaborator target:} differential geometry (symmetric spaces,
bounded domains); renormalization group in causal set theory; Wyler
programme~\citep{Wyler1969}.


\paragraph{Derivation 4: Lean proofs for baryon conservation and chirality.}
Two $\sim$50-line Lean~4 targets: (a) baryon conservation as LLC conservation
of $N_2$ spatial winding numbers, extending the existing
\texttt{RA\_AQFT\_Proofs} architecture; (b) the chirality theorem
(Proposition~\ref{prop:chirality}) formalised as a combinatorial DAG result
showing $N_3/N_4$ winding reversal requires a backward-pointing edge.
\textbf{Collaborator target:} Lean~4/Mathlib; causal set combinatorics.

\begin{remark}
The programme's history has been consistent: derivations that appeared
structural open problems have yielded to RA-native insights once stated
precisely. Charge quantization, topological confinement, parity violation,
the Koide formula, and force ranges all yielded in this paper alone. The
same expectation applies to the four remaining targets.
\end{remark}

% ── 14. Why RA is Better Positioned than LQG or String Theory ─────────────────
\section{RA vs.\ Prior Approaches}
\label{sec:comparison}

\paragraph{Loop Quantum Gravity.} LQG derives geometry successfully but imports
the Standard Model matter content as a separate input. It offers no derivation
of why the specific gauge group $U(1) \times SU(2) \times SU(3)$ or the specific
particle spectrum arises. The Bilson-Thompson braid model~\citep{BilsonThompson2005}
embedded in LQG spin networks is suggestive, but it classifies quantum numbers
without explaining why those braids are stable or how they connect to gravity.

\paragraph{String Theory.} String theory derives a rich particle spectrum but
requires 10 or 11 spacetime dimensions and a landscape of $\sim 10^{500}$
vacua~\citep{Susskind2003}. No principle selects the Standard Model vacuum.
The dimensions are extra structure added to the theory; in RA, 4 dimensions
are uniquely selected by BDG + Lovelock and require no additional input.

\paragraph{The Bilson-Thompson model.} Bilson-Thompson's topological preon
model~\citep{BilsonThompson2005} maps the first-generation Standard Model
onto braided ribbon networks with elegant results. RA is in the same general
direction, but approaches from the opposite end: instead of classifying braids
and asking which match the Standard Model, RA asks what BDG-stable patterns
the causal graph dynamics naturally produce. The BDG sign criterion provides
the stability selection rule that the Bilson-Thompson model lacks. Whether the
stable RA patterns correspond to Bilson-Thompson braids is an open question
(Derivation~1); if they do, it provides a physical derivation of why those
particular braids are the stable ones.

\paragraph{RA's unique position.} RA derives 4 dimensions (BDG + Lovelock),
three non-gravitational forces (BDG degrees of freedom in 4D), three generations
(three BDG excitation levels), no fourth generation (no $N_5$ in 4D BDG), and
no supersymmetric partners (SUSY algebra not generated by BDG dynamics) ---
all from the single ontological commitment that actualization events are
primitive. No extra dimensions are required, no vacuum landscape exists, and
the connection between gravity and the gauge forces is not a separate
unification step but an immediate consequence of the same BDG action that
already gives gravity.

% ── 15. Conclusion ───────────────────────────────────────────────────────────
\section{Conclusion}
\label{sec:conclusion}

The Relational Actualism programme began with a single commitment: that
irreversible actualization events are the primitive physical reality. From this,
it has derived quantum mechanics (exact at the discrete level), special
relativity ($c = l_P/t_P$, $E = \gamma mc^2$ from integer counting), general
relativity (inevitable by Lovelock), dark matter phenomenology (topological
dimensional reduction), and the origin of life (Erd\H{o}s-R\'{e}nyi percolation as the Causal Firewall)~\citep{ErdosRenyi1960}.

This paper extends the same logic to the Standard Model. The BDG local
topology of a 4D causal graph has exactly four integer-valued degrees of
freedom; one is gravity; three are the three gauge forces. The particle
spectrum is the catalogue of BDG-stable persistent patterns; spin comes from
self-similarity periodicity; three generations come from three BDG excitation
levels. Supersymmetry is not generated and is not needed. No fourth generation
exists because no $N_5$ exists in the 4D BDG action.

These are structural results, not phenomenological fits. Four open derivations
are precisely scoped; the expectation, given the programme's history, is that
each will yield to an insight already latent in the BDG and Poisson-CSG
machinery. The question driving the next generation of this work is exact:
is the complete set of $\SBDG > 0$-stable patterns in a 4D BDG-filtered
Poisson-CSG precisely the first-generation Standard Model particle spectrum?

If yes, Relational Actualism will have derived the Standard Model from the
single fact that some things happen and others do not.

% ── Bibliography ─────────────────────────────────────────────────────────────

\section*{Declarations}
\textbf{Funding:} Not applicable.\quad
\textbf{Competing interests:} The author declares no competing interests.

\textbf{Ethics statement:} This article does not contain any studies with human participants or animals performed by the author.

\textbf{AI use disclosure:} Claude (Anthropic, model \texttt{claude-sonnet-4-6}) was used in the preparation of this manuscript, including drafting and revision of text, mathematical derivation assistance, \LaTeX{} compilation, and Lean~4 proof development. The author is solely responsible for all scientific content, mathematical claims, and conclusions. Gemini (Google DeepMind, Gemini~2.0~Pro) was used for supplementary literature checking and cross-verification of symbolic calculations.

\textbf{Data availability:} All Lean~4 proof files and Python verification scripts are publicly available at \href{https://github.com/jsandeman/Relational-Actualism}{github.com/jsandeman/Relational-Actualism}.

\textbf{Preprint availability:} \href{https://doi.org/10.5281/zenodo.19229335}{doi.org/10.5281/zenodo.19229335}.

\bibliographystyle{plainnat}
\smallskip\noindent
\textbf{Data availability.}
All Lean~4 proof files and Python verification scripts supporting this
article are openly available at
\url{https://github.com/jsandeman/Relational-Actualism}.

\begin{thebibliography}{99}

\bibitem[Sandeman(2026BDG)]{Sandeman2026BDG}
J.~F.~Sandeman.
\newblock Two Standard Model Coupling Constants from the Unique 4-Dimensional
  BDG Integers.
\newblock \textit{Preprint}, 2026.
\newblock \href{https://doi.org/10.5281/zenodo.19324790}{doi.org/10.5281/zenodo.19324790}.

\bibitem[Gibbons \& Solodukhin(2007)]{Gibbons2007}
G.~W.~Gibbons and S.~N.~Solodukhin.
\newblock The geometry of small causal diamonds.
\newblock \textit{Physics Letters B}, \textbf{649}:317--324, 2007.
\newblock \href{https://doi.org/10.1016/j.physletb.2007.03.068}{doi:10.1016/j.physletb.2007.03.068}.

\bibitem[Sandeman(2026RACL)]{Sandeman2026RACL}
J.~F. Sandeman.
\newblock {Relational Actualism: The Einstein Field Equations as the Unique
  Macroscopic Limit of the Actualization Graph (RACL)}.
\newblock \textit{Preprint}, 2026.
\newblock \href{https://doi.org/10.5281/zenodo.19270020}{doi.org/10.5281/zenodo.19270020}.


\bibitem[Sandeman(2026RATM)]{Sandeman2026RATM}
J.~F. Sandeman.
\newblock {Relational Actualism: The Topology of Matter (RATM)}.
\newblock \textit{Preprint}, 2026.
\newblock \href{https://doi.org/10.5281/zenodo.19265651}{doi.org/10.5281/zenodo.19265651}.

\bibitem[Sandeman(2026RAEB)]{Sandeman2026RAEB}
J.~F. Sandeman.
\newblock {Relational Actualism: The Engine of Becoming (RAEB)}.
\newblock \textit{Preprint}, 2026.
\newblock \href{https://doi.org/10.5281/zenodo.19198120}{doi.org/10.5281/zenodo.19198120}.


\bibitem[Wyler(1969)]{Wyler1969}
A.~Wyler.
\newblock On the conformal groups in the theory of relativity and their
  unitary representations.
\newblock \emph{Comptes Rendus Acad.~Sci.~Paris}, 269:743--745, 1969.

\bibitem[Benincasa \& Dowker(2010)]{Benincasa2010}
D.~M.~T.~Benincasa and F.~Dowker.
\newblock The scalar curvature of a causal set.
\newblock \emph{Physical Review Letters}, 104(18):181301, 2010.
\newblock \href{https://doi.org/10.1103/PhysRevLett.104.181301}{doi:10.1103/PhysRevLett.104.181301}.

\bibitem[Bilson-Thompson(2005)]{BilsonThompson2005}
S.~O.~Bilson-Thompson.
\newblock A topological model of composite preons.
\newblock \emph{arXiv preprint}, 2005.
\newblock \href{https://arxiv.org/abs/hep-ph/0503213}{arXiv:hep-ph/0503213}.

\bibitem[Erd\H{o}s \& R\'{e}nyi(1960)]{ErdosRenyi1960}
P.~Erd\H{o}s and A.~R\'{e}nyi.
\newblock On the evolution of random graphs.
\newblock \emph{Publications of the Mathematical Institute of the Hungarian
  Academy of Sciences}, 5:17--61, 1960.

\bibitem[Lovelock(1971)]{Lovelock1971}
D.~Lovelock.
\newblock The Einstein tensor and its generalizations.
\newblock \emph{Journal of Mathematical Physics}, 12(3):498--501, 1971.
\newblock \href{https://doi.org/10.1063/1.1665613}{doi:10.1063/1.1665613}.

\bibitem[Sandeman(2026a)]{Sandeman2026RAQM}
J.~F.~Sandeman.
\newblock Relational Actualism and Quantum Mechanics ({RAQM}).
\newblock \textit{Foundations of Physics} (Under review), 2026.
\newblock \href{https://doi.org/10.5281/zenodo.19174380}{doi.org/10.5281/zenodo.19174380}.

\bibitem[Sandeman(2026b)]{Sandeman2026RAGC}
J.~F.~Sandeman.
\newblock Relational Actualism: Gravity and Cosmology ({RAGC}).
\newblock Preprint, 2026.
\newblock \href{https://doi.org/10.5281/zenodo.19197999}{doi.org/10.5281/zenodo.19197999}.



\bibitem[Sandeman(2026d)]{Sandeman2026RAQI}
J.~F.~Sandeman.
\newblock Relational Actualism: Implications for Quantum Information ({RAQI}).
\newblock Preprint, 2026.
\newblock \href{https://doi.org/10.5281/zenodo.19198285}{doi.org/10.5281/zenodo.19198285}.

\bibitem[Sandeman(2026e)]{Sandeman2026RADM}
J.~F.~Sandeman.
\newblock Relational Actualism: Dark Matter as Actualization-Density Topology ({RADM}).
\newblock Preprint, 2026.
\newblock \href{https://doi.org/10.5281/zenodo.19198513}{doi.org/10.5281/zenodo.19198513}.
\bibitem[Sandeman(2026f)]{Sandeman2026Foundation}J.~F.~Sandeman.\newblock Relational Actualism: A Foundation Paper ({Foundation}).\newblock Preprint, 2026.\newblock \href{https://doi.org/10.5281/zenodo.19229438}{doi.org/10.5281/zenodo.19229438}


\bibitem[Glaser(2014)]{Glaser2014}
L.~Glaser.
\newblock A closed form expression for the causal set d'Alembertian.
\newblock \emph{Journal of Mathematical Physics}, 55(9):092301, 2014.
\newblock \href{https://doi.org/10.1063/1.4895604}{doi:10.1063/1.4895604}.
\newblock \href{https://arxiv.org/abs/1311.1701}{arXiv:1311.1701}.

\bibitem[Hawking(1975)]{Hawking1975}
S.~W.~Hawking.
\newblock Particle creation by black holes.
\newblock \emph{Communications in Mathematical Physics}, 43(3):199--220, 1975.
\newblock \href{https://doi.org/10.1007/BF02345020}{doi:10.1007/BF02345020}.

\bibitem[Page(1993)]{Page1993}
D.~N.~Page.
\newblock Information in black hole radiation.
\newblock \emph{Physical Review Letters}, 71(23):3743--3746, 1993.
\newblock \href{https://doi.org/10.1103/PhysRevLett.71.3743}{doi:10.1103/PhysRevLett.71.3743}.

\bibitem[Particle Data Group(2022)]{PDG2022}
R.~L.~Workman et~al.
\newblock Review of particle physics.
\newblock \emph{Progress of Theoretical and Experimental Physics},
  2022(8):083C01, 2022.
\newblock \href{https://doi.org/10.1093/ptep/ptac097}{doi:10.1093/ptep/ptac097}.

\bibitem[Peccei \& Quinn(1977)]{Peccei1977}
R.~D.~Peccei and H.~R.~Quinn.
\newblock CP conservation in the presence of pseudoparticles.
\newblock \emph{Physical Review Letters}, 38(25):1440--1443, 1977.
\newblock \href{https://doi.org/10.1103/PhysRevLett.38.1440}{doi:10.1103/PhysRevLett.38.1440}.

\bibitem[Sakharov(1967)]{Sakharov1967}
A.~D.~Sakharov.
\newblock Violation of CP invariance, C asymmetry, and baryon asymmetry of the universe.
\newblock \emph{JETP Letters}, 5:24--27, 1967.

\bibitem[Susskind(1993)]{Susskind1993}
L.~Susskind, L.~Thorlacius, and J.~Uglum.
\newblock The stretched horizon and black hole complementarity.
\newblock \emph{Physical Review D}, 48(8):3743--3761, 1993.
\newblock \href{https://doi.org/10.1103/PhysRevD.48.3743}{doi:10.1103/PhysRevD.48.3743}.

\bibitem[Susskind(2003)]{Susskind2003}
L.~Susskind.
\newblock The anthropic landscape of string theory.
\newblock \emph{arXiv preprint}, 2003.
\newblock \href{https://arxiv.org/abs/hep-th/0302219}{arXiv:hep-th/0302219}.

\bibitem[Sandeman(2026g)]{Sandeman2026RAHC}
J.~F.~Sandeman.
\newblock Algorithmic Causal Dynamics: A Unified Framework for Emergence
  and Complexity in Relational Actualism ({RAHC}).
\newblock Preprint, 2026.
\newblock \href{https://doi.org/10.5281/zenodo.19198171}{doi.org/10.5281/zenodo.19198171}.

\bibitem[Sandeman(2026h)]{Sandeman2026RACI}
J.~F.~Sandeman.
\newblock Relational Actualism: Complexity and Intelligence ({RACI}).
\newblock Preprint, 2026.
\newblock \href{https://doi.org/10.5281/zenodo.19198224}{doi.org/10.5281/zenodo.19198224}.


\bibitem[Sandeman(2026RACF)]{Sandeman2026RACF}
J.~F.~Sandeman.
\newblock The Causal Firewall: Substrate-Independent Physical Constraints
  on the Distribution of Life in the Universe (RACF).
\newblock \textit{Int.\ J.\ Astrobiol.} (submitted), 2026.
\newblock \href{https://doi.org/10.5281/zenodo.19319374}%
  {doi:10.5281/zenodo.19319374}.

\end{thebibliography}

\end{document}
